\section{Introduction}


\subsection{}
\label{1.1}
Let 
%
\begin{equation}
E(\mathbf{w}) 
= 
\frac{1}{2} 
\sum_{n = 1}^{N} 
\left( 
	y(x_n, 
	\mathbf{w}) 
	- 
	t_n 
\right) ^ 2.
\end{equation}
%
Setting the derivative to zero gives
%
\begin{equation}
\mathbf{0} 
= 
\sum_{n = 1}^{N} 
\frac{\partial y(x_n, \mathbf{w})}{\partial \mathbf{w}} 
\left( 
	y(x_n, \mathbf{w}) 
	- 
	t_n 
\right).
\end{equation}
%
Substituting 
%
\begin{equation}
y(x_n, \mathbf{w}) 
= 
\sum_{j = 0}^{M} 
w_j 
x_n ^ j
\end{equation}
%
gives
%
\begin{equation}
0 
= 
\sum_{n = 1}^{N} 
x_n ^ i 
\left( 
	\sum_{j = 0}^{M} 
	w_j 
	x_n ^ j 
	- 
	t_n 
\right).
\end{equation}
%
Therefore,
%
\begin{equation}
\argmin_{w_j}
E(\mathbf{w})
=
\left \{
	w_j
	\left.
	\right \vert
	\sum_{j = 0}^{M} 
	A_{ij} 
	w_j 
	= 
	T_i	
\right \},
\end{equation}
%
where
%
\begin{equation}
\begin{aligned}
A_{ij} 
&= 
\sum_{n = 1}^{N} 
x_n^{i + j}, \\
T_i 
&= 
\sum_{n = 1}^{N} 
x_n ^i 
t_n.
\end{aligned}
\end{equation}


\subsection{}
\label{1.2}
Let
%
\begin{equation}
\tilde{E}(\mathbf{w}) 
= 
\frac{1}{2} 
\sum_{n = 1}^{N} 
\left( 
	y(x_n, \mathbf{w}) 
	- 
	t_n 
\right) ^ 2 
+ 
\frac{\lambda}{2} 
\lVert \mathbf{w} \rVert ^ 2.
\end{equation}
%
Setting the derivative to zero gives
%
\begin{equation}
\mathbf{0} 
= 
\sum_{n = 1}^{N} 
\frac{\partial y(x_n, \mathbf{w})}{\partial \mathbf{w}} 
\left( 
	y(x_n, \mathbf{w}) 
	- 
	t_n 
\right) 
+ 
\lambda 
\mathbf{w}.
\end{equation}
%
Substituting 
%
\begin{equation}
y(x_n, \mathbf{w}) 
= 
\sum_{j = 0}^{M} 
w_j 
x_n ^ j
\end{equation}
%
gives
%
\begin{equation}
0 
= 
\sum_{n = 1}^{N} 
x_n ^ i 
\left( 
	\sum_{j = 0}^{M} 
	w_j 
	x_n ^ j 
	- 
	t_n 
\right) 
+ 
\lambda 
w_i.
\end{equation}
%
Therefore,
%
\begin{equation}
\argmin_{w_j}
E(\mathbf{w})
=
\left \{
	w_j
	\left.
	\right \vert
	\sum_{j = 0}^{M} 
	\tilde{A}_{ij} 
	w_j 
	= 
	T_i	
\right \},
\end{equation}
%
where
%
\begin{equation}
\begin{aligned}
\tilde{A}_{ij} 
&= 
\sum_{n = 1}^{N} 
x_n ^ {i + j} 
+ 
\lambda 
\delta_{ij}, \\
T_i 
&= 
\sum_{n = 1}^{N} 
x_n ^ i 
t_n.
\end{aligned}
\end{equation}


\subsection{}
\label{1.3}
Let $a$, $o$ and $l$ be the events where an apple, orange and lime are selected respectively.
The probability that an apple is selected is given by
%
\begin{equation}
p(a) 
= 
p(a | r) 
p(r) 
+ 
p(a | b) 
p(b) 
+ 
p(a | g) 
p(g).
\end{equation}
%
Substituting 
$p(a | r) = \frac{3}{10}$, 
$p(r) = \frac{1}{5}$, 
$p(a | g) = \frac{1}{2}$, 
$p(r) = \frac{1}{5}$, 
$p(a | g) = \frac{3}{10}$ 
and 
$p(g) = \frac{3}{5}$ 
gives
%
\begin{equation}
p(a) 
= 
\frac{17}{50}.
\end{equation}
%

If an orange is selected, the probability that it came from the geen box is given by
%
\begin{equation}
p(g | o) = \frac{p(g, o)}{p(o)}.
\end{equation}
%
Here,
%
\begin{equation}
\begin{aligned}
p(g, o) &= p(o | g) p(g), \\
p(o) & = p(o | r) p(r) + p(o | b) p(b) + p(o | g) p(g).
\end{aligned}
\end{equation}
%
Substituting 
$p(o | r) = \frac{2}{5}$, 
$p(r) = \frac{1}{5}$, 
$p(o | b) = \frac{1}{2}$, 
$p(b) = \frac{1}{5}$, 
$p(o | g) = \frac{3}{10}$ 
and 
$p(g) = \frac{3}{5}$ 
gives 
$p(g, o) = \frac{9}{50}$ 
and 
$p(o) = \frac{9}{25}$.
%
Therefore,
\begin{equation}
p(g | o) = \frac{1}{2}.
\end{equation}


\subsection{}
\label{1.4}
Let 
%
\begin{equation}
x 
= 
g(y)
\end{equation}
%
and $\hat{x}$ and $\hat{y}$ be the locations of the maximum of $p_x(x)$ and $p_y(y)$ respcetively.
Let us assume that there exists $\epsilon > 0$ such that $g'(y) \neq 0$ for $\left| y - \hat{y} \right| < \epsilon$.
Then, Taking the derivative of the transoformation
%
\begin{equation}
p_y(y) 
= 
p_x 
\left( 
	g(y) 
\right) 
\abs{g'(y)}
\end{equation}
%
and substituting $y = \hat{y}$ gives
%
\begin{equation}
0 
= 
g'(\hat{y}) 
p_x' 
\left( 
	g 
	\left( 
		\hat{y} 
	\right) 
\right) 
+ 
p_x 
\left( 
	g 
	\left( 
		\hat{y} 
	\right) 
\right) 
g'' 
\left( 
	\hat{y} 
\right).
\end{equation}
%
Therefore, in general,
%
\begin{equation}
\hat{x} 
\neq 
g 
\left( 
	\hat{y} 
\right).
\end{equation}
%

Here, let us assume that 
%
\begin{equation}
g(y) 
= 
a y 
+ 
b.
\end{equation}
%
Then, Taking the derivative of the transformation and substituting $y = \hat{y}$ gives
%
\begin{equation}
0 
= 
p_x' 
\left( 
	g 
	\left( 
		\hat{y} 
	\right) 
\right).
\end{equation}
%
Therefore, 
%
\begin{equation}
\hat{x} 
= 
g 
\left( 
	\hat{y} 
\right).
\end{equation}
%


\subsection{}
\label{1.5}
By the definition, 
%
\begin{equation}
\Var f(x) 
= 
\E 
\left( 
	f(x) 
	- 
	\E f(x)
\right) ^ 2.
\end{equation}
%
The right hand side can be written as 
%
\begin{equation}
\E 
\left( 
	\left( 
		f(x) 
	\right) ^ 2 
	- 
	2 
	f(x) 
	\E f(x) 
	+ 
	\left( 
		\E f(x) 
	\right) ^ 2 
\right) 
= 
\E 
\left( 
	f(x) 
\right) ^ 2 
- 
\left( 
	\E f(x) 
\right) ^ 2.
\end{equation}
%
Therefore, 
%
\begin{equation}
\Var f(x) 
= 
\E 
\left( 
	f(x) 
\right) ^ 2 
- 
\left( 
	\E f(x) 
\right) ^ 2.
\end{equation}
%


\subsection{}
\label{1.6}
By the definition,
%
\begin{equation}
\cov (x, y) 
= 
\E 
\left( 
	\left( 
		x 
		- 
		\E x 
	\right) 
	\left( 
		y 
		- 
		\E y
	\right) 
\right).
\end{equation}
%
The right hand side can be written as
%
\begin{equation}
\E x y 
- 
\E 
\left( 
	x 
	\E y 
\right) 
- 
\E 
\left( 
	y \E x 
\right) 
+ 
\E 
\left( 
	\E x 
	\E y 
\right) 
= 
\E x y 
- 
\E x 
\E y.
\end{equation}
%
The right hand side can be written as
%
\begin{equation}
\int 
x 
y 
p(x, y) 
dx 
dy 
- 
\int 
x 
p(x) 
dx 
\int 
y 
p(y) 
dy.
\end{equation}
%
If $x$ and $y$ are independent, by the definition,
%
\begin{equation}
f(x, y) 
= 
f(x) 
f(y).
\end{equation}
%
Then,
%
\begin{equation}
\int 
x 
y 
p(x, y) 
dx 
dy 
= 
\int 
p(x) 
dx 
\int 
p(y) 
dy.
\end{equation}
%
Therefore,
%
\begin{equation}
\cov (x, y) 
= 
0.
\end{equation}
%


\subsection{}
\label{1.7}
Let
%
\begin{equation}
I 
= 
\int_{-\infty}^{\infty} 
\exp 
\left( 
	-
	\frac{1}{2 \sigma ^ 2} 
	x ^ 2 
\right) 
dx.
\end{equation}
%
Then
%
\begin{equation}
I ^ 2 
= 
\int_{-\infty}^{\infty} 
\int_{-\infty}^{\infty} 
\exp 
\left( 
	-
	\frac{1}{2 \sigma ^ 2} 
	\left( 
		x ^ 2 
		+ 
		y ^ 2 
	\right) 
\right) 
dx 
dy.
\end{equation}
%
By the transformation from Cartesian coordinates $(x, y)$ to polar coordinates $(r, \theta)$, the right hand side can be written as
%
\begin{equation}
\int_{0}^{\infty} 
\int_{0}^{2 \pi} 
\exp 
\left( 
	- 
	\frac{1}{2 \sigma ^ 2} 
	r ^ 2 
\right) 
\begin{vmatrix}
\cos \theta & - r \sin \theta \\
\sin \theta &  r \cos \theta \\
\end{vmatrix}
dr 
d\theta
= 
2 \pi 
\int_{0}^{\infty} 
\exp 
\left( 
	- 
	\frac{1}{2 \sigma ^ 2} 
	r ^ 2 
\right) 
r 
dr.
\end{equation}
%
By the transformation $s = \frac{r}{\sigma}$, the right hand side can be written as
%
\begin{equation}
2 \pi \sigma ^ 2 
\int_{0}^{\infty} 
\exp 
\left( 
	- 
	\frac{1}{2} 
	s ^ 2 
\right)
s 
ds 
= 
2 \pi \sigma ^ 2 
\left[ 
	- 
	\exp 
	\left( 
		- 
		\frac{1}{2} 
		s ^ 2 
	\right) 
\right]_{0}^{\infty}.
\end{equation}
%
Therefore, 
%
\begin{equation}
I 
= 
\left( 
	2 \pi \sigma ^ 2 
\right) ^ \frac{1}{2}. 
\end{equation}
%

By the definition,
%
\begin{equation}
\mathcal{N} 
\left( 
	x 
	| 
	\mu, 
	\sigma ^ 2 
\right) 
= 
\left( 
	2 \pi \sigma ^ 2 
\right) ^ {- \frac{1}{2}} 
\exp 
\left( 
	- 
	\frac{1}{2 \sigma ^ 2} 
	(x - \mu) ^ 2 
\right).
\end{equation}
%
Then
%
\begin{equation}
\int_{- \infty}^{\infty} 
\mathcal{N} 
\left( 
	x 
	| 
	\mu, 
	\sigma ^ 2 
\right) 
dx 
= 
\left( 
	2 \pi \sigma ^ 2 
\right) ^ {- \frac{1}{2}} 
\int_{-\infty}^{\infty} 
\exp 
\left( 
	-
	\frac{1}{2 \sigma ^ 2} 
	(x - \mu) ^ 2 
\right) 
dx.
\end{equation}
%
By the transformation $t = x - \mu$, the right hand side can be written as 
%
\begin{equation}
\left( 
	2 \pi \sigma ^ 2 
\right) ^ {- \frac{1}{2}} 
\int_{-\infty}^{\infty} 
\exp 
\left( 
	-
	\frac{1}{2 \sigma ^ 2} 
	t ^ 2 
\right) dt 
= 
\left( 
	2 \pi \sigma ^ 2 
\right) ^ {- \frac{1}{2}} 
I.
\end{equation}
%
Therefore,
%
\begin{equation}
\int_{- \infty}^{\infty} 
\mathcal{N} 
\left( 
	x 
	| 
	\mu, 
	\sigma ^ 2 
\right) 
dx 
= 
1.
\end{equation}
%


\subsection{}
\label{1.8}
Let $x$ be a variable such that
%
\begin{equation}
p(x) 
= 
\mathcal{N} (x | \mu, \sigma ^ 2).
\end{equation}
%
Then
%
\begin{equation}
\E x 
= 
\int_{-\infty}^{\infty} 
x 
\mathcal{N} \left( x | \mu, \sigma ^ 2 \right) 
dx.
\end{equation}
%
By the definition, the right hand side can be written as
%
\begin{equation}
\left( 
	2 \pi \sigma ^ 2 
\right) ^ {-\frac{1}{2}} 
\int_{-\infty}^{\infty} 
x 
\exp 
\left( 
	-
	\frac{1}{2 \sigma ^ 2} 
	(x - \mu) ^ 2 
\right) 
dx.
\end{equation}
%
By the transformation $y = x - \mu$, it can be written as 
%
\begin{equation}
\left( 
	2 \pi \sigma ^ 2 
\right) ^ {-\frac{1}{2}} 
\int_{-\infty}^{\infty} 
(y + \mu) 
\exp 
\left( 
	-
	\frac{1}{2 \sigma ^ 2}
	y ^ 2 
\right) 
dy.
\end{equation}
%
Since 
%
\begin{equation}
\left( 
	2 \pi \sigma ^ 2 
\right) ^ {-\frac{1}{2}} 
\int_{-\infty}^{\infty} 
y 
\exp 
\left( 
	-
	\frac{1}{2 \sigma ^ 2}
	y ^ 2 
\right) 
dy 
= 
0,
\end{equation}
%
and
%
\begin{equation}
\left( 
	2 \pi \sigma ^ 2 
\right) ^ {-\frac{1}{2}} 
\int_{-\infty}^{\infty} 
	\mu 
	\exp 
	\left( 
		-
		\frac{1}{2 \sigma ^ 2}
		y ^ 2 
	\right) 
dy 
= 
\mu 
\int_{- \infty}^{\infty} 
\mathcal{N} 
\left( 
	y 
	| 
	\mu, 
	\sigma ^ 2 
\right) 
dy,
\end{equation}
%
we have
%
\begin{equation}
\E x 
= 
\mu.
\end{equation}
%

By the definition,  
%
\begin{equation}
\int_{- \infty}^{\infty} 
\mathcal{N} 
\left( 
	x 
	| 
	\mu, 
	\sigma ^ 2 
\right) 
dx 
= 
1
\end{equation}
%
can be written as
%
\begin{equation}
\left( 
	2 \pi \sigma ^ 2 
\right) ^ {-\frac{1}{2}} 
\int_{-\infty}^{\infty} 
\exp 
\left( 
	- 
	\frac{1}{2 \sigma ^ 2} 
	(x - \mu) ^ 2 
\right) 
dx 
= 
1.
\end{equation}
%
Taking the derivative with respect to $\sigma ^ 2$ gives 
%
\begin{equation}
\begin{aligned}
\left( 
	2 \pi 
\right) ^ {-\frac{1}{2}} 
\left( 
	- 
	\frac{1}{2} 
\right) 
\left( 
	\sigma ^ 2 
\right) ^ {- \frac{3}{2}} 
\int_{-\infty}^{\infty} 
\exp 
\left( 
	-
	\frac{1}{2 \sigma ^ 2} 
	(x - \mu) ^ 2 
\right) dx \\
+ 
\left( 
	2 \pi \sigma ^ 2 
\right) ^ {-\frac{1}{2}} 
\int_{-\infty}^{\infty} 
\frac{1}{2} 
\left( 
	\sigma ^ 2 
\right) ^ {- 2} 
(x - \mu) ^ 2 
\exp 
\left( 
	- 
	\frac{1}{2 \sigma ^ 2} 
	(x - \mu) ^ 2 
\right) 
dx 
= 
0.
\end{aligned}
\end{equation}
%
The left hand side can be written as
%
\begin{equation}
\begin{aligned}
- 
\frac{1}{2} 
\left( 
	\sigma ^ 2 
\right) ^ {- 1} 
\int_{- \infty}^{\infty} 
\mathcal{N} 
\left( 
	x 
	| 
	\mu, 
	\sigma ^ 2 
\right) 
dx 
+ 
\frac{1}{2} 
\left( 
	\sigma ^ 2 
\right) ^ {- 2} 
\int_{- \infty} ^ {\infty} 
(x - \mu) ^ 2 
\mathcal{N} 
\left( 
	x 
	| 
	\mu, 
	\sigma ^ 2 
\right) dx \\
= 
- 
\frac{1}{2} 
\left( 
	\sigma ^ 2 
\right) ^ {- 1} 
+ 
\frac{1}{2} 
\left( 
	\sigma ^ 2 
\right) ^ {- 2} 
\Var x.
\end{aligned}
\end{equation}
%
Therefore,
%
\begin{equation}
\Var x 
= 
\sigma ^ 2.
\end{equation}
%


\subsection{}
\label{1.9}
Let
%
\begin{equation}
\mathcal{N} 
\left( 
	x 
	| 
	\mu, 
	\sigma ^ 2 
\right) 
= 
\left( 
	2 \pi \sigma ^ 2 
\right) ^ {- \frac{1}{2}} 
\exp 
\left( 
	- 
	\frac{1}{2 \sigma ^ 2} 
	(x - \mu) ^ 2 
\right).
\end{equation}
%
Setting its derivative with respect to $x$ to zero gives
%
\begin{equation}
0 
= 
\left( 
	2 \pi \sigma ^ 2 
\right) ^ {- \frac{1}{2}} 
\left( 
	- 
	\frac{1}{\sigma ^ 2} 
	(x - \mu) 
\right) 
\exp 
\left( 
	- 
	\frac{1}{2 \sigma ^ 2} 
	(x - \mu) ^ 2 
\right).
\end{equation}
%
Therefore, the mode is given by $\mu$.

Similarly, let
%
\begin{equation}
\mathcal{N} 
\left( 
	\mathbf{x} 
	| 
	\bm{\mu}, 
	\bm{\Sigma} 
\right) 
= 
\left( 
	2 \pi 
\right) ^ {- \frac{D}{2}} 
(\det \bm{\Sigma}) ^ {- \frac{1}{2}} 
\exp 
\left( 
	- 
	\frac{1}{2} 
	(\mathbf{x} - \bm{\mu}) ^ \intercal 
	\bm{\Sigma} ^ {- 1} 
	(\mathbf{x} - \bm{\mu}) 
\right).
\end{equation}
%
Setting its derivative with respect to $\mathbf{x}$ to zero gives
%
\begin{equation}
\mathbf{0} 
= 
- 
\left( 
	2 \pi 
\right) ^ {- \frac{D}{2}} 
(\det \bm{\Sigma}) ^ {- \frac{1}{2}} 
\left( 
	\bm{\Sigma} ^ {- 1} 
	+ 
	\left( 
		\bm{\Sigma} ^ {- 1}
	\right) ^ \intercal 
\right) 
(\mathbf{x} - \bm{\mu}) 
\exp 
\left( 
	- 
	\frac{1}{2} 
	(\mathbf{x} - \bm{\mu}) ^ \intercal 
	\bm{\Sigma} ^ {- 1} 
	(\mathbf{x} - \bm{\mu}) 
\right).
\end{equation}
%
Therefore, the mode is given by $\bm{\mu}$.


\subsection{}
\label{1.10}
By the definition,
%
\begin{equation}
\E (x + y) 
= 
\int 
\int 
(x + y) 
p(x, y) 
dx 
dy.
\end{equation}
%
The right hand side can be written as 
%
\begin{equation}
\int 
x 
\left( 
	\int 
	p(x, y) 
	dy 
\right) 
dx 
+ 
\int 
y 
\left( 
	\int 
	p(x, y) 
	dx 
\right) 
dy 
= 
\int 
x 
p(x) 
dx 
+ 
\int 
y 
p(y) 
dy.
\end{equation}
%
By the definition, the right hand side can be written as 
%
\begin{equation}
\E x 
+ 
\E y.
\end{equation}
%
Therefore, 
%
\begin{equation}
\E (x + y) 
= 
\E x 
+ 
\E y.
\end{equation}
%

Similarly, by the definition,
%
\begin{equation}
\Var (x + y) 
= 
\E 
\left( 
	x 
	+ 
	y 
	- 
	\E (x + y) 
\right) ^ 2
\end{equation}
%
By the result above and the definition, the right hand side can be written as 
%
\begin{equation}
\begin{aligned}
\E 
\left( 
	x 
	- 
	\E x 
\right) ^ 2 
+ 
2 
\E 
\left( 
	\left( 
		x 
		- 
		\E x 
	\right) 
	\left( 
		y 
		- 
		\E y 
	\right) 
\right) 
+ 
\E 
\left( 
	y 
	- 
	\E y 
\right) ^ 2 \\
= 
\Var x 
+ 
2 \cov (x, y) 
+ 
\Var y.
\end{aligned}
\end{equation}
%
If $x$ and $y$ are independent, then 
%
\begin{equation}
\cov (x, y) 
= 
0,
\end{equation}

%
by \ref{1.6}. Therefore,
%
\begin{equation}
\Var (x + y) 
= 
\Var x 
+ 
\Var y.
\end{equation}
%


\subsection{}
\label{1.11}
Let $\mathbf{x}$ be a set of $N$ variables such that
%
\begin{equation}
\ln p 
\left( 
	\mathbf{x} 
	| 
	\mu, 
	\sigma ^ 2 
\right) 
= 
- 
\frac{N}{2} 
\ln 
\left( 
	2 \pi \sigma ^ 2 
\right) 
- 
\frac{1}{2 \sigma ^ 2} 
\sum_{n = 1}^{N} 
(x_n - \mu) ^ 2.
\end{equation}
%
To maximise it with respect to $\mu$ and $\sigma ^ 2$, setting the partial derivatives to zero gives
%
\begin{equation}
\begin{aligned}
0 
&= 
\frac{1}{\sigma ^ 2} 
\sum_{n = 1}^{N} 
(x_n - \mu), \\
0 
&= 
- 
\frac{N}{2 \sigma ^ 2} 
+ 
\frac{1}{2 \left( \sigma ^ 2 \right) ^ 2} 
\sum_{n = 1}^{N} 
(x_n - \mu) ^ 2.
\end{aligned}
\end{equation}
%
Therefore,
%
\begin{equation}
\begin{aligned}
\mu_{\rm ML} 
&= 
\frac{1}{N} 
\sum_{n = 1}^{N} 
x_n, \\
\sigma_{\rm ML} ^ 2 
&= 
\frac{1}{N} 
\sum_{n = 1}^{N} 
(x_n - \mu_{\rm ML}) ^ 2.
\end{aligned}
\end{equation}
%


\subsection{}
\label{1.12}
Let $x_m$ and $x_n$ be independent variables. 
Then
%
\begin{equation}
\E x_m x_n 
= 
\E x_m 
\E x_n.
\end{equation}
%
If they are samples from the Gaussian distribution with mean $\mu$ and variance $\sigma ^ 2$, the right hand side is given by $\mu ^ 2$.
On the other hand, by the definition, 
%
\begin{equation}
\E x_n ^ 2 
=  
\Var x_n 
+ 
\left( 
	\E x_n 
\right) ^ 2.
\end{equation}
%
If $x_n$ is a sample from the Gaussian distribution with mean $\mu$ and variance $\sigma ^ 2$, the right hand side is given by $\sigma ^ 2 + \mu ^ 2$.
Therefore,
%
\begin{equation}
\E x_m x_n 
= 
\mu ^ 2 
+ 
\delta_{mn}
\sigma ^ 2.
\end{equation}
%

Here, since 
%
\begin{equation}
\mu_{\rm ML} 
= 
\frac{1}{N} 
\sum_{n = 1}^{N} 
x_n,
\end{equation}
%
we have
%
\begin{equation}
\E \mu_{\rm ML} 
= 
\frac{1}{N} 
\sum_{n = 1}^{N} 
\E x_n.
\end{equation}
%
Therefore, 
%
\begin{equation}
\E \mu_{\rm ML} 
= 
\mu.
\end{equation}
%
Similarly, since 
%
\begin{equation}
\sigma_{\rm ML}^2 
= 
\frac{1}{N} 
\sum_{n = 1}^{N} 
\left( 
	x_n 
	- 
	\mu_{\rm ML} 
\right) ^ 2,
\end{equation}
%
we have 
%
\begin{equation}
\E \sigma_{\rm ML}^2 
= 
\frac{1}{N} 
\sum_{n = 1}^{N} 
\E 
\left( 
	x_n 
	- 
	\mu_{\rm ML} 
\right) ^ 2.
\end{equation}
%
The right hand side can be writen as 
%
\begin{equation}
\frac{1}{N} 
\sum_{n = 1}^{N} 
\E 
\left( 
	x_n ^ 2 
	- 
	2 \mu_{\rm ML} x_n 
	+ 
	\mu_{\rm ML} ^ 2 
\right) 
= 
\frac{1}{N} 
\sum_{n = 1}^{N} 
\E x_n ^ 2 
- 
\frac{2}{N} 
\E 
\left( 
	\mu_{\rm ML} 
	\left( 
		\sum_{n = 1}^{N} 
		x_n 
	\right) 
\right) 
+ 
\E \mu_{\rm ML} ^ 2.
\end{equation}
%
The first term of the right hand side can be written as 
%
\begin{equation}
\frac{1}{N}
\sum_{n = 1}^{N} 
\left( 
	\mu ^ 2 
	+ 
	\sigma ^ 2 
\right) 
= 
\mu ^ 2 
+ 
\sigma ^ 2,
\end{equation}
%
while the second and third terms can be writen as
%
\begin{equation}
- 
2 
\E \mu_{\rm ML} ^ 2 
+ 
\E \mu_{\rm ML} ^ 2 
= 
- 
\E \mu_{\rm ML} ^ 2.
\end{equation}
%
Here, 
%
\begin{equation}
\E \mu_{\rm ML} ^ 2 
= 
\E 
\left( 
	\frac{1}{N} 
	\sum_{n = 1}^{N} x_n 
\right) ^ 2.
\end{equation}
%
The right hand side can be written as 
%
\begin{equation}
\frac{1}{N ^ 2} 
\sum_{n = 1}^{N} 
\E x_n ^ 2 
+ 
\frac{2}{N ^ 2} 
\sum_{1 \leq m < n \leq N} 
\E x_m x_n 
= 
\frac{1}{N} 
\left(
	\mu ^ 2 
	+ 
	\sigma ^ 2
\right) 
+ 
\frac{N - 1}{N} 
\mu ^ 2.
\end{equation}
%
Therefore,
%
\begin{equation}
\E \mu_{\rm ML} ^ 2 
= 
\mu ^ 2 
+ 
\frac{1}{N} \sigma ^ 2.
\end{equation}
%
Thus,
%
\begin{equation}
\E \sigma_{\rm ML}^2 
= 
\frac{N - 1}{N} \sigma ^ 2.
\end{equation}
%


\subsection{}
\label{1.13}
Let $\{ x_n \}$ be a set of variables whose mean is $\mu$ and variance is $\sigma ^ 2$.
Then
%
\begin{equation}
\E 
\left( 
	\frac{1}{N} 
	\sum_{n = 1}^{N} 
	\left(
		x_n 
		- 
		\mu 
	\right) ^ 2 
\right) 
= 
\frac{1}{N} 
\sum_{n = 1}^{N} 
\E 
\left( 
	x_n 
	- 
	\mu 
\right) ^ 2.
\end{equation}
%
The right hand side can be writen as 
%
\begin{equation}
\frac{1}{N} 
\sum_{n = 1}^{N} 
\E 
\left( 
	x_n ^ 2 
	- 
	2 \mu x_n 
	+ 
	\mu ^ 2 
\right) 
= 
\frac{1}{N} 
\sum_{n = 1}^{N} 
\E x_n ^ 2 
- 
\frac{2 \mu}{N}  
\sum_{n = 1}^{N} 
\E x_n 
+ 
\mu ^ 2.
\end{equation}
%
The first term of the right hand side can be written as 
%
\begin{equation}
\frac{1}{N} 
\sum_{n = 1}^{N} 
\left( 
	\mu ^ 2 
	+ 
	\sigma ^ 2 
\right) 
= 
\mu ^ 2 
+ 
\sigma ^ 2,
\end{equation}
%
while the second term can be writen as
%
\begin{equation}
- 
\frac{2 \mu}{N} 
\sum_{n = 1}^{N} 
\mu 
= 
- 
2 
\mu ^ 2.
\end{equation}
%
Therefore,
%
\begin{equation}
\E 
\left( 
	\frac{1}{N} 
	\sum_{n = 1}^{N} 
	\left( 
		x_n 
		- 
		\mu 
	\right) ^ 2 
\right) 
= 
\sigma ^2.
\end{equation}
%


\subsection{}
\label{1.14}
Let
%
\begin{equation}
\begin{aligned}
w_{ij}^{\rm S} 
&= 
\frac{1}{2} 
(w_{ij} 
+ 
w_{ji}), \\
w_{ij}^{\rm A} 
&= 
\frac{1}{2} 
(w_{ij}
 -
  w_{ji}).
\end{aligned}
\end{equation}
%
Then
%
\begin{equation}
\begin{aligned}
w_{ij} 
&= 
w_{ij}^{\rm S} 
+ 
w_{ij}^{\rm A}, \\
w_{ij}^{\rm S} 
&=
w_{ji}^{\rm S}, \\
w_{ij}^{\rm A} 
&= 
- 
w_{ji}^{\rm A}.
\end{aligned}
\end{equation}
%
Here,
\begin{equation}
\sum_{i = 1}^{D} 
\sum_{j = 1}^{D} 
w_{ij}^{\rm A} 
x_i x_j 
= 
\frac{1}{2} 
\sum_{i = 1}^{D} 
\sum_{j = 1}^{D} 
(w_{ij} 
- 
w_{ji}) 
x_i x_j.
\end{equation}
%
The right hand side can be written as
%
\begin{equation}
\frac{1}{2} 
\left( 
	\sum_{i = 1}^{D} 
	\sum_{j = 1}^{D} 
	w_{ij} x_i x_j 
	- 
	\sum_{i = 1}^{D} 
	\sum_{j = 1}^{D} 
	w_{ji} x_i x_j 
\right) 
= 
0.
\end{equation}
%
Therefore, 
%
\begin{equation}
\sum_{i = 1}^{D} 
\sum_{j = 1}^{D} 
w_{ij}^{\rm A} x_i x_j 
= 
0.
\end{equation}
%

Additionally,
%
\begin{equation}
\sum_{i = 1}^{D} 
\sum_{j = 1}^{D} 
w_{ij} x_i x_j 
= 
\sum_{i = 1}^{D} 
\sum_{j = 1}^{D} 
\left( 
	w_{ij}^{\rm S} 
	+ 
	w_{ij}^{\rm A} 
\right) 
x_i x_j.
\end{equation}
%
The right hand side can be written as
%
\begin{equation}
\sum_{i = 1}^{D} 
\sum_{j = 1}^{D} 
w_{ij}^{\rm S} x_i x_j 
+ 
\sum_{i = 1}^{D} 
\sum_{j = 1}^{D} 
w_{ij}^{\rm A} x_i x_j 
= 
\sum_{i = 1}^{D} 
\sum_{j = 1}^{D} 
w_{ij}^{\rm S} x_i x_j,
\end{equation}
%
where the result above is used. 
Therefore,
%
\begin{equation}
\sum_{i = 1}^{D} 
\sum_{j = 1}^{D} 
w_{ij} x_i x_j 
= 
\sum_{i = 1}^{D} 
\sum_{j = 1}^{D} 
w_{ij}^{\rm S} x_i x_j.
\end{equation}
%

Finally, since the matrix $w_{ij}^{\rm S}$ is a $D \times D$ symmetric matrix, its number of independent parameters is $\frac{D (D + 1)}{2}$.


\subsection{}
\label{1.15}
Let $n(D, M)$ be the number of independent parameters of a polynomial in $D$ dimensions and $M$ orders.
Then
%
\begin{equation}
n(1, M)
=
n(1, M - 1)
=
1.
\end{equation}
%
Let us assume that
%
\begin{equation}
n(D, M)
=
\sum_{i = 1}^{D}
n(i, M - 1).
\end{equation}
%
The independent terms of a polynomial in $D + 1$ dimensions and $M$ orders can be split into 1. the ones of a polynomial in $D$ dimensions and $M$ orders and 2. the ones generated by multiplying the ones in $D + 1$ dimensions and $M$ orders by the $D + 1$th variable.
Therefore,
%
\begin{equation}
n(D + 1, M)
=
n(D, M)
+
n(D + 1, M - 1).
\end{equation}
%
Thus,
%
\begin{equation}
n(D + 1, M)
=
\sum_{i = 1}^{D + 1}
n(i, M - 1).
\end{equation}
%
Hence, the assumption is proved by induction on $D$.

Additionally,
%
\begin{equation}
\sum_{i = 1}^{1}
\frac{(i + M - 2)!}{(i - 1)! (M - 1)!}
=
1.
\end{equation}
%
Let us assume that
%
\begin{equation}
\sum_{i = 1}^{D}
\frac{(i + M - 2)!}{(i - 1)! (M - 1)!}
=
\frac{(D + M - 1)!}{(D - 1)! M!}.
\end{equation}
%
Then
%
\begin{equation}
\sum_{i = 1}^{D + 1}
\frac{(i + M - 2)!}{(i - 1)! (M - 1)!}
=
\frac{(D + M - 1)!}{(D - 1)! M!}
+
\frac{(D + M - 1)!}{D! (M - 1)!}.
\end{equation}
%
The right hand side can be written as
%
\begin{equation}
\frac{D (D + M - 1)! + M (D + M - 1)!}{D! M!}
=
\frac{(D + M)!}{D! M!}.
\end{equation}
%
Therefore, the assumption is proved by induction on $D$.

Finally, by \ref{1.14},
%
\begin{equation}
n(D, 2)
=
\frac{D (D + 1)}{2}.
\end{equation}
%
Let us assume that
%
\begin{equation}
n(D, M)
=
\frac{(D + M - 1)!}{(D - 1)! M!}.
\end{equation}
%
Then, by the result above,
%
\begin{equation}
n(D, M + 1)
=
\sum_{i = 1}^{D}
n(i, M).
\end{equation}
%
By the assumption and result above, the right hand side can be written as
%
\begin{equation}
\sum_{i = 1}^{D}
\frac{(i + M - 1)!}{(i - 1)! M!}
=
\frac{(D + M)!}{(D - 1)! (M + 1)!}.
\end{equation}
%
Therefore, the assumption is proved by induction on $M$.


\subsection{}
\label{1.16}
Let $N(D, M)$ be the number of independent parameters in all of the terms up to and including the ones of $D$ dimensions and $M$ orders.
Then, by \ref{1.15},
%
\begin{equation}
N(D, M)
=
\sum_{m = 0}^{M}
n(D, m),
\end{equation}
%
where
%
\begin{equation}
n(D, m)
=
\frac{(D + m - 1)!}{(D - 1)! m!}.
\end{equation}
%

Additionally,
%
\begin{equation}
N(D, 0)
=
1.
\end{equation}
%
Let us assume that
%
\begin{equation}
\sum_{m = 0}^{M}
n(D, m)
=
\frac{(D + M)!}{D! M!}.
\end{equation}
%
Then
%
\begin{equation}
\sum_{m = 0}^{M + 1}
n(D, m)
=
\frac{(D + M)!}{D! M!}
+
\frac{(D + M)!}{(D - 1)! (M + 1)!}.
\end{equation}
%
The right hand side can be written as
%
\begin{equation}
\frac{(M + 1) (D + M)! + D (D + M)!}{D! (M + 1)!}
=
\frac{(D + M + 1)!}{D! (M + 1)!}.
\end{equation}
%
Therefore, the assumption is proved by induction on $M$.
Thus,
%
\begin{equation}
N(D, M)
=
\frac{(D + M)!}{D! M!}.
\end{equation}
%

Additionally, by the approximation
%
\begin{equation}
n!
\simeq
n ^ n
\exp(- n),
\end{equation}
%
the right hand side can be approximated as
%
\begin{equation}
\frac
{(D + M) ^ {D + M} 
\exp 
\left(
	- (D + M)
\right)}
{D ^ D \exp (- D) M ^ M \exp(- M)}
=
\frac{(D + M) ^ {D + M}}{D ^ D M ^ M}.
\end{equation}
%
The right hand side can be written as
%
\begin{equation}
D ^ M 
\left(
	1
	+
	\frac{M}{D}
\right) ^ D
\left(
	\frac{1}{M}
	+
	\frac{1}{D}
\right) ^ M
=
M ^ D 
\left(
	1
	+
	\frac{D}{M}
\right) ^ M
\left(
	\frac{1}{D}
	+
	\frac{1}{M}
\right) ^ D.
\end{equation}
%
Therefore, $N(D, M)$ can be approximated as $D ^ M$ for $D \gg M$ and as $M ^ D$ for $M \gg D$.

Finally, by the result above, 
%
\begin{equation}
\begin{aligned}
N(10, 3)
&=
286, \\
N(100, 3)
&=
176851, \\
N(1000, 3)
&=
167668501.
\end{aligned}
\end{equation}
%


\subsection{}
\label{1.17}
Let
%
\begin{equation}
\Gamma(x) 
= 
\int_{0}^{\infty} 
u ^ {x - 1} 
\exp(- u) 
du.
\end{equation}
%
Then
%
\begin{equation}
\Gamma(x + 1) 
= 
\int_{0}^{\infty} 
u ^ {x} 
\exp(- u) 
du.
\end{equation}
%
The right hand side can be written as 
%
\begin{equation}
\left[ 
	- u ^ x 
	\exp(- u) 
\right]_{u = 0}^{u = \infty} 
+ 
\int_{0}^{\infty} 
x 
u ^ {x - 1} 
\exp(- u) 
du 
= 
x 
\Gamma(x).
\end{equation}
%
Therefore,
%
\begin{equation}
\Gamma( x + 1) 
= 
x \Gamma(x).
\end{equation}
%

Since
%
\begin{equation}
\Gamma(1) 
= 
\int_{0}^{1} 
\exp(- u) 
du,
\end{equation}
%
and the right hand side can be written as $1$,
%
\begin{equation}
\Gamma(1) 
= 
0!.
\end{equation}
%
For a positive integer $x$, let us assume that
%
\begin{equation}
\Gamma(x) 
=
 (x - 1)!.
\end{equation}
%
Then,
%
\begin{equation}
\Gamma(x + 1) 
= 
x 
\Gamma(x),
\end{equation}
%
where the right hand side can be written as $x!$.
Therefore,
%
\begin{equation}
\Gamma(x + 1) = x!.
\end{equation}
%
Thus, the assumption is proved by induction on $x$.


\subsection{}
\label{1.18}
Let us consider the transformation from Cartesian to polar coordinates
%
\begin{equation}
\prod_{i = 1}^{D} 
\int_{- \infty}^{\infty} 
\exp 
\left( 
	- x_i ^ 2 
\right) 
dx_i 
= 
S_D 
\int_{0}^{\infty} 
\exp 
\left( 
	- r ^ 2 
\right) 
r ^ {D - 1} 
dr,
\end{equation}
%
where $S_D$ is the surface area of a sphere of unit raidus in $D$ dimensions. 
By \ref{1.7}, the left hand side can be written as $\pi ^ \frac{D}{2}$. By the transformation $s = r ^ 2$, the right hand side can be written as
%
\begin{equation}
\frac{S_D}{2} 
\int_{0}^{\infty} 
\exp(- s) 
s ^ \frac{D - 1}{2} 
s ^ {- \frac{1}{2}} 
ds 
= 
\frac{S_D}{2} 
\Gamma 
\left( 
	\frac{D}{2} 
\right).
\end{equation}
%
Therefore,
%
\begin{equation}
S_D 
= 
\frac{2 \pi ^ \frac{D}{2}}
{\Gamma 
\left( 
	\frac{D}{2} 
\right)}.
\end{equation}
%

Additionally, the volume of the sphere can can be written as
%
\begin{equation}
V_D 
= 
S_D 
\int_{0}^{1} 
r^{D - 1} 
dr.
\end{equation}
%
The right hand side can be written as 
%
\begin{equation}
S_D 
\left[ 
	\frac{r ^ D}{D} 
\right]_{r = 0}^{r = 1} 
= 
\frac{S_D}{D}.
\end{equation}
%
Therefore, 
%
\begin{equation}
V_D 
= 
\frac{S_D}{D}.
\end{equation}
%

Finally, the results above reduce to
%
\begin{equation}
\begin{aligned}
S_2 
&= 
\frac{2 \pi}{\Gamma(1)}, \\
V_2 
&= 
\frac{S_2}{2}. 
\end{aligned}
\end{equation}
%
Therefore, 
%
\begin{equation}
\begin{aligned}
S_2 
&= 
2 \pi, \\
V_2 
&= 
\pi. 
\end{aligned}
\end{equation}
%
Similarly,
%
\begin{equation}
\begin{aligned}
S_3 
&= 
\frac{2 \pi ^ \frac{3}{2}}
{\Gamma 
\left( 
	\frac{3}{2} 
\right)}, \\
V_3 
&= 
\frac{S_3}{3}. 
\end{aligned}
\end{equation}
%
Therefore, 
%
\begin{equation}
\begin{aligned}
S_3 
&= 
4 \pi, \\
V_3 
&= 
\frac{4}{3} \pi. 
\end{aligned}
\end{equation}
%


\subsection{}
\label{1.19}
The volume of a cube of side 2 in $D$ dimensions is $2 ^ D$. 
Therefore, the ratio of the volume of the cocentric sphere of radius 1 divided by the volume of the cube is given by
%
\begin{equation}
\frac{V_D}{2 ^ D} 
= 
\frac{\pi ^ \frac{D}{2}}{D 2 ^ {D - 1} 
\Gamma 
\left( 
	\frac{D}{2} 
\right)},
\end{equation}
%
by \ref{1.18}.

Additionally, by Stering's formula
%
\begin{equation}
\Gamma(x + 1) 
\simeq 
(2 \pi) ^ \frac{1}{2} 
\exp(- x) 
x ^ \frac{x + 1}{2},
\end{equation}
%
the ratio can be approximated as
%
\begin{equation}
\frac{V_D}{2 ^ D} 
\simeq 
\frac
{\pi ^ \frac{D}{2}}
{D 2 ^ {D - 1} 
(2 \pi) ^ \frac{1}{2} 
\exp 
\left( 
	1 
	- 
	\frac{D}{2} 
\right) 
\left( 
	\frac{D}{2} 
	- 
	1 
\right) ^ \frac{D}{4}}.
\end{equation}
%
The right hand side can be written as 
%
\begin{equation}
\frac{1}{2 e (2 \pi) ^ \frac{1}{2}} 
\frac{1}{D} 
\left( 
	\frac{e ^ 2 \pi ^ 2}{ 8 D - 16} 
\right) ^ \frac{D}{4}.
\end{equation}
%
Therefore, the ratio goes to zero as $D \to \infty$.

Finally, the ratio of the distance from the center of the cube to one of the corners divided by the perpendicular distance to one of the sides is given by
%
\begin{equation}
\frac{\sqrt{\sum_{i = 1}^{D} 1 ^ 2}}{1} 
= 
\sqrt{D}.
\end{equation}
%
Therefore, the ration goes to $\infty$ as $D \to \infty$.


\subsection{}
\label{1.20}
For a vector $\mathbf{x}$ in $D$ dimensions, let 
%
\begin{equation}
p(\mathbf{x}) 
= 
(2 \pi \sigma ^ 2) ^ {- \frac{D}{2}} 
\exp 
\left( 
	- 
	\frac{\norm{\mathbf{x}} ^ 2}{2 \sigma ^2} 
\right).
\end{equation}
%
Then
%
\begin{equation}
\int_{r \leq \norm{\mathbf{x}} \leq r + \epsilon} 
p(\mathbf{x}) 
d\mathbf{x} 
= 
\int_{r}^{r + \epsilon} 
\int (2 \pi \sigma ^ 2) ^ {- \frac{D}{2}} 
\exp 
\left( 
	- 
	\frac{{r'} ^ 2}{2 \sigma ^2} 
\right) 
J 
dr' 
d\bm{\phi},
\end{equation}
%
where $\bm{\phi}$ is the vector of the angular components of the polar corrdinate and $J$ is the Jacobian of the transformation from the Cartesian to polar coordinate.
For a sufficiently small $\epsilon$, the right hand side can be approximated as
%
\begin{equation}
\begin{aligned}
(2 \pi \sigma ^ 2) ^ {- \frac{D}{2}} 
\exp 
\left( 
	- 
	\frac{r ^ 2}{2 \sigma ^2} 
\right) 
\int_{r}^{r + \epsilon} 
\int 
J 
dr' 
d\bm{\phi} \\
= 
(2 \pi \sigma ^ 2) ^ {- \frac{D}{2}} 
\exp 
\left( 
	- 
	\frac{r ^ 2}{2 \sigma ^2} 
\right) 
\int_{r \leq \norm{\mathbf{x}} \leq r + \epsilon} 
d\mathbf{x}.
\end{aligned}
\end{equation}
%
Therefore,
%
\begin{equation}
\int_{r \leq \norm{\mathbf{x}} \leq r + \epsilon} 
p(\mathbf{x}) 
d\mathbf{x} 
\simeq 
p(r) 
\epsilon,
\end{equation}
%
where
%
\begin{equation}
p(r) 
= 
(2 \pi \sigma ^ 2) ^ {- \frac{D}{2}} 
S_D 
r ^ {D - 1} 
\exp 
\left( 
	- 
	\frac{r ^ 2}{2 \sigma ^2} 
\right),
\end{equation}
%
and $S_D$ is the surface area of a unit sphere in $D$ dimensions.

Additionally, setting the derivative of $p(r)$ to zero gives
%
\begin{equation}
0 
= 
(2 \pi \sigma ^ 2) ^ {- \frac{D}{2}} 
S_D 
\left( 
	(D - 1) r ^ {D - 2} 
	- 
	\frac{r ^ D}{\sigma ^ 2} 
\right) 
\exp 
\left( 
	- 
	\frac{r ^ 2}{2 \sigma ^2} 
\right).
\end{equation}
%
Therefore, $p(r)$ is maximised at a sigle stationary point
%
\begin{equation}
\hat{r} 
= 
\sqrt{D - 1} 
\sigma.
\end{equation}
%

Additionally, by the expression of $p(r)$ above,
%
\begin{equation}
\frac{p 
\left( 
	\hat{r} 
	+ 
	\epsilon 
\right)}
{p(\hat{r})} 
= 
\left( 
	\frac{\hat{r} + \epsilon}{\hat{r}} 
\right) ^ {D - 1} 
\exp 
\left( 
	- 
	\frac{2 \hat{r} \epsilon + \epsilon ^ 2}{2 \sigma ^ 2} 
\right).
\end{equation}
%
Using the expression of $\hat{r}$ above, the right hand side can be written as
%
\begin{equation}
\begin{aligned}
\exp 
\left( 
	(D - 1) 
	\ln 
	\left( 
		1 
		+ 
		\frac{\epsilon}{\hat{r}} 
	\right) 
	- 
	\frac{2 \hat{r} \epsilon + \epsilon ^ 2}{2 \sigma ^ 2} 
\right) \\
= \exp 
\left( 
	\frac{\hat{r} ^ 2}{\sigma ^ 2} 
	\ln 
	\left( 
		1 
		+ 
		\frac{\epsilon}{\hat{r}} 
	\right) 
	- 
	\frac{2 \hat{r} \epsilon + \epsilon ^ 2}{2 \sigma ^ 2} 
\right).
\end{aligned}
\end{equation}
%
By the Taylor series
%
\begin{equation}
\ln (1 + x) 
= 
x 
- 
\frac{1}{2} x ^ 2 
+ 
o \left( x ^ 3 \right),
\end{equation}
%
the right hand side can be approximated as
%
\begin{equation}
\exp 
\left( 
	\frac{\hat{r} ^ 2}{\sigma ^ 2} 
	\left( 
		\frac{\epsilon}{\hat{r}} 
		- 
		\frac{\epsilon ^ 2}{2 {\hat{r}} ^ 2} 
	\right) 
	- 
	\frac{2 \hat{r} \epsilon + \epsilon ^ 2}{2 \sigma ^ 2} 
\right) 
= 
\exp 
\left( 
	- 
	\frac{\epsilon ^ 2}{\sigma ^ 2} 
\right).
\end{equation}
%
Therefore,
%
\begin{equation}
p 
\left( 
	\hat{r} 
	+ 
	\epsilon 
\right) 
\simeq 
p 
\left( 
	\hat{r} 
\right) 
\exp 
\left( 
	- 
	\frac{\epsilon ^ 2}{\sigma ^ 2} 
\right).
\end{equation}
%

Finally, let a vector of length $\hat{r}$ be $\hat{\mathbf{r}}$.
Then, by the definition of $p(\mathbf{x})$,
%
\begin{equation}
\frac
{p(\mathbf{0})}
{p 
\left( 
	\hat{\mathbf{r}} 
\right)} 
= 
\exp 
\left( 
	\frac{\hat{r} ^ 2}{2 \sigma ^2} 
\right).
\end{equation}
%
Substituting the expression of $\hat{r}$ above, the right hand side can be written as $\exp \left( \frac{D - 1}{2} \right)$.
Therefore,
%
\begin{equation}
\frac
{p(\mathbf{0})}
{p 
\left( 
	\hat{\mathbf{r}} 
\right)} 
= 
\exp 
\left( 
	\frac{D - 1}{2} 
\right).
\end{equation}
%


\subsection{}
\label{1.21}
If $0 \leq a \leq b$, then
%
\begin{equation}
0 
\leq 
a (b - a).
\end{equation}
%
Therefore,
%
\begin{equation}
a 
\leq 
(ab) ^ \frac{1}{2}.
\end{equation}
%

For a two-class classification problem of $\mathbf{x}$, let the classes be $\mathcal{C}_1$ and $\mathcal{C}_2$ and let the decision regions be $\mathcal{R}_1$ and $\mathcal{R}_2$.
Let us choose the decision regions to minimise the probability of misclassification.
Then,
%
\begin{equation}
p(\mathbf{x}, \mathcal{C}_1) 
> 
p(\mathbf{x}, \mathcal{C}_2) 
\Rightarrow 
\mathbf{x} \in \mathcal{C}_1, 
\end{equation}
%
and
%
\begin{equation}
p(\mathbf{x}, \mathcal{C}_2) 
> 
p(\mathbf{x}, \mathcal{C}_1) 
\Rightarrow \mathbf{x} 
\in 
\mathcal{C}_2.
\end{equation}
%
Then, using the inequality above,
%
\begin{equation}
\int_{\mathcal{R}_1} 
p(\mathbf{x}, \mathcal{C}_2) 
d\mathbf{x} 
\leq 
\int_{\mathcal{R}_1} 
\left( 
	p(\mathbf{x}, \mathcal{C}_1) 
	p(\mathbf{x}, \mathcal{C}_2) 
\right) ^ \frac{1}{2} 
d\mathbf{x},
\end{equation}
%
and
%
\begin{equation}
\int_{\mathcal{R}_2} 
p(\mathbf{x}, \mathcal{C}_1) 
d\mathbf{x} 
\leq 
\int_{\mathcal{R}_2} 
\left( 
	p(\mathbf{x}, \mathcal{C}_1) 
	p(\mathbf{x}, \mathcal{C}_2) 
\right) ^ \frac{1}{2} 
d\mathbf{x}.
\end{equation}
%
Therefore,
%
\begin{equation}
\int_{\mathcal{R}_1} 
p(\mathbf{x}, \mathcal{C}_2) 
d\mathbf{x} 
+ 
\int_{\mathcal{R}_2} 
p(\mathbf{x}, \mathcal{C}_1) 
d\mathbf{x} 
\leq 
\int 
\left( 
	p(\mathbf{x}, \mathcal{C}_1) 
	p(\mathbf{x}, \mathcal{C}_2) 
\right) ^ \frac{1}{2} 
d\mathbf{x}.
\end{equation}
%


\subsection{}
\label{1.22}
Let
%
\begin{equation}
\E L 
= 
\sum_{k} 
\sum_{j} 
\int_{\mathcal{R}_j} L_{kj} 
p(\mathbf{x}, \mathcal{C}_k) 
d\mathbf{x}.
\end{equation}
%
If
%
\begin{equation}
L_{kj} 
= 
1 
- 
\delta_{kj},
\end{equation}
%
then the right hand side can be written as
%
\begin{equation}
\sum_{k} \sum_{j} 
\int_{\mathcal{R}_j} 
\left( 
	p(\mathbf{x}, \mathcal{C}_k) 
	- 
	p(\mathbf{x}, \mathcal{C}_j) 
\right)  
d\mathbf{x} 
= 
\sum_{j} 
\int_{\mathcal{R}_j} 
\left( 
	\sum_{k} 
	p(\mathbf{x}, \mathcal{C}_k) 
	- 
	p(\mathbf{x}, \mathcal{C}_j) 
\right)  
d\mathbf{x}.
\end{equation}
%
The right hand side can be written as
%
\begin{equation}
\sum_{j} 
\int_{\mathcal{R}_j} 
\left( 
	p(\mathbf{x}) 
	- 
	p(\mathbf{x}, \mathcal{C}_j) 
\right)  
d\mathbf{x} 
= 
1 
- 
\sum_{j} 
\int_{\mathcal{R}_j} 
p(\mathbf{x}, \mathcal{C}_j) 
d\mathbf{x}.
\end{equation}
%
Therefore,
%
\begin{equation}
\E L 
= 
1 
- 
\sum_{j} 
\int_{\mathcal{R}_j}
p(\mathcal{C}_j | \mathbf{x}) 
p(\mathbf{x}) 
d\mathbf{x}.
\end{equation}
%
Thus, minimising $\E L$ reduces to choosing the criterion to maximise the posterior probatility $p(\mathcal{C}_j | \mathbf{x})$.


\subsection{}
\label{1.23}
Let
%
\begin{equation}
\E L 
= 
\sum_{k} 
\sum_{j} 
\int_{\mathcal{R}_j} 
L_{kj} 
p(\mathbf{x}, \mathcal{C}_k) 
d\mathbf{x}.
\end{equation}
%
The right hand side can be written as
%
\begin{equation}
\sum_{j} 
\int_{\mathcal{R}_j} 
\sum_{k} 
L_{kj} 
p(\mathbf{x}, \mathcal{C}_k) 
d\mathbf{x} 
= 
\sum_{j} 
\int_{\mathcal{R}_j} 
\left( 
	\sum_{k} 
	L_{kj} 
	p(\mathcal{C}_k | \mathbf{x}) 
\right) 
p(\mathbf{x}) 
d\mathbf{x}.
\end{equation}
%
Therefore,
%
\begin{equation}
\E L 
= 
\sum_{j} 
\int_{\mathcal{R}_j} 
\left( 
	\sum_{k} 
	L_{kj} 
	p(\mathcal{C}_k | \mathbf{x}) 
\right) 
p(\mathbf{x}) 
d\mathbf{x}.
\end{equation}
%
Thus, minimising $\E L$ reduces to choosing to minimise $\sum_{k} L_{kj} p(\mathcal{C}_k | \mathbf{x})$.


\subsection{(Incomplete)}
\label{1.24}
Let
%
\begin{equation}
\E L 
= 
\sum_{k} 
\sum_{j} 
\int_{\mathcal{R}_j} 
L_{kj} 
p(\mathbf{x}, \mathcal{C}_k) 
d\mathbf{x}
+
\lambda
\int_{\forall k p(\mathcal{C}_k | \mathbf{x}) < \theta}
p(\mathbf{x})
d\mathbf{x}.
\end{equation}
%


\subsection{}
\label{1.25}
Let
%
\begin{equation}
\E L(\mathbf{t}, \mathbf{y} (\mathbf{x})) 
= 
\int 
\int 
\norm{\mathbf{y}(\mathbf{x}) - \mathbf{t}} ^ 2 
p(\mathbf{x}, \mathbf{t}) 
d\mathbf{x} 
d\mathbf{t}.
\end{equation}
%
Setting the derivative with respect to $\mathbf{y}(\mathbf{x})$ to zero gives
%
\begin{equation}
\mathbf{0}
= 
2 
\int 
\left( 
	\mathbf{y} (\mathbf{x}) - \mathbf{t} 
\right) 
p(\mathbf{x}, \mathbf{t}) 
d\mathbf{t}.
\end{equation}
%
The integral of the right hand side can be written as
%
\begin{equation}
\mathbf{y}(\mathbf{x}) 
\int 
p(\mathbf{x}, \mathbf{t}) 
d\mathbf{t} 
- 
\int 
\mathbf{t} p(\mathbf{x}, \mathbf{t}) 
d\mathbf{t} 
=
\mathbf{y}(\mathbf{x})
p(\mathbf{x})
-
p(\mathbf{x})
\int
\mathbf{t}
p(\mathbf{t} | \mathbf{x})
d\mathbf{t}.
\end{equation}
%
The integral in the second term of the right hand side can be written as $\E_\mathbf{t} (\mathbf{t} | \mathbf{x})$.
Therefore, the right hand side can be written as
%
\begin{equation}
\mathbf{0}
=
p(\mathbf{x})
\left(
	\mathbf{y}(\mathbf{x}) 
	- 
	\E_\mathbf{t} (\mathbf{t} | \mathbf{x})
\right).
\end{equation}
%
Thus,
%
\begin{equation}
\argmin_{\mathbf{y}(\mathbf{x})}
\E L(\mathbf{t}, \mathbf{y} (\mathbf{x}))
= 
\E_\mathbf{t} (\mathbf{t} | \mathbf{x}).
\end{equation}
%

Finally, for a single target variable $t$, it reduces to 
%
\begin{equation}
\argmin_{\mathbf{y}(\mathbf{x})}
\E L(\mathbf{t}, \mathbf{y} (\mathbf{x}))
= 
\E_t (t | \mathbf{x}).
\end{equation}
%


\subsection{}
\label{1.26}
Let
%
\begin{equation}
\E L(\mathbf{t}, \mathbf{y} (\mathbf{x})) 
= 
\int 
\int 
\norm{\mathbf{y}(\mathbf{x}) - \mathbf{t}} ^ 2 
p(\mathbf{x}, \mathbf{t}) 
d\mathbf{x} 
d\mathbf{t}.
\end{equation}
%
The right hand side can be written as
%
\begin{equation}
\begin{aligned}
&
\int 
\int \norm{\mathbf{y}(\mathbf{x}) - \E_\mathbf{t} (\mathbf{t} | \mathbf{x}) + \E_\mathbf{t} (\mathbf{t} | \mathbf{x}) - \mathbf{t}} ^ 2 
p(\mathbf{x}, \mathbf{t}) 
d\mathbf{x} 
d\mathbf{t} \\ 
=& 
\int 
\int \norm{\mathbf{y}(\mathbf{x}) - \E_\mathbf{t} (\mathbf{t} | \mathbf{x})} ^ 2 p(\mathbf{x}, \mathbf{t}) 
d\mathbf{x} 
d\mathbf{t} \\
&+ 
2 
\int 
\int 
\left( 
	\mathbf{y}(\mathbf{x}) 
	- 
	\E_\mathbf{t} (\mathbf{t} | \mathbf{x}) 
\right) ^ \intercal 
\left( 
	\E_\mathbf{t} (\mathbf{t} | \mathbf{x}) 
	- 
	\mathbf{t} 
\right) 
p(\mathbf{x}, \mathbf{t}) 
d\mathbf{x} 
d\mathbf{t} \\
&+ 
\int 
\int 
\norm{\E_\mathbf{t} (\mathbf{t} | \mathbf{x}) - \mathbf{t}} ^ 2 
p(\mathbf{x}, \mathbf{t}) 
d\mathbf{x} 
d\mathbf{t}.
\end{aligned}
\end{equation}
%
Let us look at each term of the right hand side.
The first term can be written as
%
\begin{equation}
\int 
\norm{\mathbf{y}(\mathbf{x}) - \E_\mathbf{t} (\mathbf{t} | \mathbf{x})} ^ 2 \left( 
	\int 
	p(\mathbf{x}, \mathbf{t}) 
	d\mathbf{t} 
\right) 
d\mathbf{x} 
= 
\int 
\norm{\mathbf{y}(\mathbf{x}) - \E_\mathbf{t} (\mathbf{t} | \mathbf{x})} ^ 2 p(\mathbf{x}) 
d\mathbf{x}.
\end{equation}
%
The second term can be written as
%
\begin{equation}
2 
\int 
\left( 
	\mathbf{y}(\mathbf{x}) 
	- 
	\E_\mathbf{t} (\mathbf{t} | \mathbf{x}) 
\right) ^ \intercal 
\left( 
	\int 
	\left( 
		\E_\mathbf{t} (\mathbf{t} | \mathbf{x}) - \mathbf{t} 
	\right) 
	p(\mathbf{t} | \mathbf{x}) 
	d\mathbf{t} 
\right) 
p(\mathbf{x}) 
d\mathbf{x}.
\end{equation}
%
Since
%
\begin{equation}
\begin{aligned}
\int 
\E_\mathbf{t} (\mathbf{t} | \mathbf{x}) 
p(\mathbf{t} | \mathbf{x}) 
d\mathbf{t} 
&= 
\E_\mathbf{t} (\mathbf{t} | \mathbf{x}) 
\frac
{\int p(\mathbf{x}, \mathbf{t}) d\mathbf{t}}
{p(\mathbf{x})}, \\
\int 
\mathbf{t} 
p(\mathbf{t} | \mathbf{x}) 
d\mathbf{t} 
&= 
\E_\mathbf{t} (\mathbf{t} | \mathbf{x}),
\end{aligned}
\end{equation}
%
the second term is zero.
The third term can be written as
%
\begin{equation}
\int 
\left( 
	\int 
	\norm{\E_\mathbf{t} (\mathbf{t} | \mathbf{x}) - \mathbf{t}} ^ 2 
	p(\mathbf{t} | \mathbf{x}) 
	d\mathbf{t} 
\right) 
p(\mathbf{x}) 
d\mathbf{x} 
= 
\int 
\Var_{\mathbf{t}} (\mathbf{t} | \mathbf{x}) 
p(\mathbf{x}) 
d\mathbf{x}.
\end{equation}
%
Therefore,
%
\begin{equation}
\E L(\mathbf{t}, \mathbf{y} (\mathbf{x})) 
= 
\int 
\norm{\mathbf{y}(\mathbf{x}) - \E_\mathbf{t} (\mathbf{t} | \mathbf{x})} ^ 2 p(\mathbf{x}) 
d\mathbf{x} 
+ 
\int 
\Var_{\mathbf{t}} (\mathbf{t} | \mathbf{x}) p(\mathbf{x}) 
d\mathbf{x}.
\end{equation}
%
Thus, 
%
\begin{equation}
\argmin_{\mathbf{y}(\mathbf{x})} 
\E L(\mathbf{t}, \mathbf{y} (\mathbf{x}))
= 
\E_\mathbf{t} (\mathbf{t} | \mathbf{x}).
\end{equation}
%.


\subsection{(Incomplete)}
\label{1.27}
Let
%
\begin{equation}
\E L_q 
= 
\int 
\int 
\abs{y(\mathbf{x}) - t} ^ q 
p(\mathbf{x}, t) 
d\mathbf{x} 
dt.
\end{equation}
%
Setting the derivative with respect to $y(\mathbf{x})$ to zero gives
%
\begin{equation}
0 
= 
q
p(\mathbf{x})
\int  
\abs{y(\mathbf{x}) - t} ^ {q - 1} 
{\rm sign} (y(\mathbf{x}) - t) 
p(t | \mathbf{x}) 
dt.
\end{equation}
%
Therefore,
%
\begin{equation}
\argmin_{y(\mathbf{x})}
\E L_q
=
\left \{
	y(\mathbf{x})
	\left.
	\right \vert
	\int  
	\abs{y(\mathbf{x}) - t} ^ {q - 1} 
	{\rm sign} (y(\mathbf{x}) - t) 
	p(t | \mathbf{x}) 
	dt
	=
	0
\right \}.
\end{equation}
%

Additionally, if $q = 1$, the integral can be written as
%
\begin{equation}
p(\mathbf{x})
\int
{\rm sign}(y(\mathbf{x}) - t)
p(t | \mathbf{x})
dt
=
p(\mathbf{x})
\left(
	\int_{y(\mathbf{x})}^{\infty} 
	p(t | \mathbf{x}) 
	dt 
	- 
	\int_{- \infty}^{y(\mathbf{x})}
	p(t | \mathbf{x}) 
	dt
\right).
\end{equation}
%
Therefore, 
%
\begin{equation}
\argmin_{y(\mathbf{x})}
\E L_1
=
\median (t | \mathbf{x}).
\end{equation}
%

Finally, 
%
\begin{equation}
\lim_{q \to 0}
\argmin_{y(\mathbf{x})}
\E L_q
=
\mode (t | \mathbf{x}) ?
\end{equation}
%


\subsection{}
\label{1.28}
Let us assume that
%
\begin{equation}
p(x, y) 
= 
p(x) 
p(y) 
\Rightarrow 
h(x, y) 
= 
h(x) 
+ 
h(y).
\end{equation}
%
Let $h(p)$ be a function to relate $h$ and $p$.
Then
%
\begin{equation}
h 
\left( 
	p ^ 2 
\right)
= 
2
h(p).
\end{equation}
%
Let us assume that, for a positive integer $n$,
%
\begin{equation}
h 
\left( 
	p ^ n 
\right) 
= 
n 
h(p).
\end{equation}
%
Then, by the first assumption,
%
\begin{equation}
h 
\left( 
	p ^ {n + 1} 
\right) 
= 
h 
\left( 
	p ^ n 
\right) 
+ 
h(p).
\end{equation}
%
Therefore,
%
\begin{equation}
h 
\left( 
	p ^ {n + 1} 
\right) 
= 
(n + 1) 
h(p).
\end{equation}
%
Thus, the second assumption is proved by induction on $n$.

Additionally, for positive integers $m$ and $n$, 
%
\begin{equation}
h 
\left( 
	p ^ n 
\right) 
= 
h 
\left(
	 p ^ {\frac{n}{m} m} 
\right).
\end{equation}
%
By the second assumption, the left hand side can be written as $n h(p)$.
By the first assumption, the right hand side can be written as $m h \left( p ^ \frac{n}{m} \right)$.
Therefore,
%
\begin{equation}
h 
\left(
	p ^ \frac{n}{m} 
\right) 
= 
\frac{n}{m} 
h(p).
\end{equation}
%

Finally, by the continuity, for a positive real number $a$,
%
\begin{equation}
h 
\left( 
	p ^a 
\right) 
= 
a 
h(p).
\end{equation}
%
Taking the derivative with respect to $a$ and substituting $a = 1$ gives
%
\begin{equation}
(p \ln p) 
h' (p) 
= 
h(p).
\end{equation}
%
Therefore,
%
\begin{equation}
\int 
\frac{h'(p)}{h(p)} 
dp 
= 
\int 
\frac{1}{p \ln p} 
dp 
+ 
\const.
\end{equation}
%
Ignorting the constants, the left hand side can be written as $\ln h(p)$ and the right hand side can be written as $\ln (\ln p)$.
Thus,
%
\begin{equation}
h(p) 
\propto 
\ln p.
\end{equation}
%


\subsection{}
\label{1.29}
Let $x$ be an $M$-state discrete random variable.
Then, by the definition,
%
\begin{equation}
{\rm H}(x) 
= 
- 
\sum_{i = 1}^{M} 
p(x_i) 
\ln p(x_i),
\end{equation}
%
where
%
\begin{equation}
\sum_{i = 1}^{M} 
p(x_i) 
= 
1.
\end{equation}
%
By Jensen's inequality,
%
\begin{equation}
\sum_{i = 1}^{M} 
p(x_i) 
\ln \frac{1}{p(x_i)} 
\leq 
\ln 
\left( 
	\sum_{i = 1}^{M} 1 
\right).
\end{equation}
Therefore,
%
\begin{equation}
{\rm H}(x) 
\leq 
\ln M.
\end{equation}
%


\subsection{}
\label{1.30}
Let 
%
\begin{equation}
\begin{aligned}
p(x) 
&= 
\mathcal{N} 
\left( 
	x 
	| 
	\mu, 
	\sigma ^2 
\right), \\
q(x) 
&= 
\mathcal{N} 
\left( 
	x 
	| 
	m, 
	s ^2 
\right).
\end{aligned}
\end{equation}
%
By the definition,
%
\begin{equation}
{\rm KL} 
\left( 
	p 
	|| 
	q 
\right) 
= 
- 
\int 
p(x) 
\ln 
\frac{q(x)}{p(x)} 
dx.
\end{equation}
%
The right hand side can be written as 
%
\begin{equation}
\begin{aligned}
&
- 
\int_{- \infty}^{\infty} 
p(x) 
\ln 
\frac{
\left( 
	2 \pi s ^ 2 
\right) ^ {- \frac{1}{2}} 
\exp 
\left( 
	- 
	\frac{(x - m) ^ 2}{2 s ^ 2} 
\right)
}
{
\left( 
	2 \pi \sigma ^ 2 
\right) ^ {- \frac{1}{2}} 
\exp 
\left( 
	- 
	\frac{(x - \mu) ^ 2}{2 \sigma ^ 2} 
\right)} 
dx \\
& 
= 
- 
\int_{- \infty}^{\infty} 
p(x) 
\left( 
	-
	\frac{1}{2} 
	\ln \frac{s ^ 2}{\sigma ^ 2} 
	- 
	\frac{(x - m) ^ 2}{2 s ^ 2} 
	+ 
	\frac{(x - \mu) ^ 2}{2 \sigma ^ 2} 
\right) 
dx.
\end{aligned}
\end{equation}
%
The right hand side can be written as
%
\begin{equation}
\ln \frac{s}{\sigma} 
\int_{- \infty}^{\infty} 
p(x) 
dx 
+ 
\frac{1}{2 s ^2} 
\int_{- \infty}^{\infty} 
(x - m) ^ 2 
p(x) 
dx 
- 
\frac{1}{2 \sigma ^2} 
\int_{- \infty}^{\infty} 
(x - \mu) ^ 2 
p(x) 
dx.
\end{equation}
%
%
The first term can be written as $\ln \frac{s}{\sigma}$.
%
The second term can be written as
%
\begin{equation}
\frac{1}{2 s ^2} 
\int_{- \infty}^{\infty} 
(x - \mu + \mu - m) ^ 2 
p(x) 
dx 
= 
\frac{\sigma ^ 2 + (\mu - m) ^ 2}{2 s ^ 2}.
\end{equation}
%
The third term can be written as $- \frac{1}{2}$.
Therefore,
%
\begin{equation}
{\rm KL} 
\left( 
	p 
	|| 
	q 
\right) 
= 
\ln \frac{s}{\sigma} 
+ 
\frac{\sigma ^ 2 + (\mu - m) ^ 2}{2 s ^ 2} 
- 
\frac{1}{2}. 
\end{equation}
%


\subsection{}
\label{1.31}
Let $\mathbf{x}$ and $\mathbf{y}$ be two variables.
Then, by the definition, 
%
\begin{equation}
\begin{aligned}
{\rm H}(\mathbf{x}) 
&= 
- 
\int
p(\mathbf{x}) 
\ln p(\mathbf{x}) 
d\mathbf{x}, \\
{\rm H}(\mathbf{y}) 
&= 
- 
\int 
p(\mathbf{y}) 
\ln p(\mathbf{y}) 
d\mathbf{y}, \\
{\rm H}(\mathbf{x}, \mathbf{y}) 
&= 
- 
\int 
\int 
p(\mathbf{x}, \mathbf{y}) 
\ln p(\mathbf{x}, \mathbf{y}) 
d\mathbf{x} 
d\mathbf{y}.
\end{aligned}
\end{equation}
%
Note that
%
\begin{equation}
\begin{aligned}
{\rm H}(\mathbf{x}) 
&= 
- 
\int 
\left( 
	\int 
	p(\mathbf{x}, \mathbf{y}) 
	d\mathbf{y} 
\right) 
\ln p(\mathbf{x}) 
d\mathbf{x}, \\
{\rm H}(\mathbf{y}) 
&= 
- 
\int 
\left( 
	\int p(\mathbf{x}, \mathbf{y}) 
	d\mathbf{x} 
\right) 
\ln p(\mathbf{y}) 
d\mathbf{y}.
\end{aligned}
\end{equation}
%
Therefore,
%
\begin{equation}
{\rm H}(\mathbf{x}) 
+ 
{\rm H}(\mathbf{y}) 
- 
{\rm H}(\mathbf{x}, \mathbf{y}) 
= 
- 
\int 
\int 
p(\mathbf{x}, \mathbf{y}) 
\ln \frac{p(\mathbf{x}) p(\mathbf{y})}{p(\mathbf{x}, \mathbf{y})} 
d\mathbf{x} 
d\mathbf{y}.
\end{equation}
%
Since
%
\begin{equation}
\int 
\int 
p(\mathbf{x}, \mathbf{y}) 
d\mathbf{x} 
d\mathbf{y} 
=
1,
\end{equation}
%
Jensen's inequality can be used to write that
%
\begin{equation}
- 
\int 
\int 
p(\mathbf{x}, \mathbf{y}) 
\ln \frac{p(\mathbf{x}) p(\mathbf{y})}{p(\mathbf{x}, \mathbf{y})} 
d\mathbf{x} d\mathbf{y} 
\geq 
- 
\ln 
\left( 
	\int 
	\int 
	p(\mathbf{x}) 
	p(\mathbf{y}) 
	d\mathbf{x} 
	d\mathbf{y} 
\right).
\end{equation}
%
The right hand side can be written as
%
\begin{equation}
- 
\ln 
\left( 
	\int 
	p(\mathbf{x}) 
	d\mathbf{x} 
	\int  
	p(\mathbf{y}) 
	d\mathbf{y} 
\right) 
= 
0.
\end{equation}
%
Thus,
%
\begin{equation}
{\rm H}(\mathbf{x}, \mathbf{y}) 
\leq 
{\rm H}(\mathbf{x}) 
+ 
{\rm H}(\mathbf{y}).
\end{equation}
%


\subsection{}
\label{1.32}
Let $\mathbf{x}$ be a vector of continuous variables and 
%
\begin{equation}
\mathbf{y} 
= 
\mathbf{A} 
\mathbf{x},
\end{equation}
%
where $\mathbf{A}$ is a nonsingular matrix.
By the definition, 
%
\begin{equation}
{\rm H}(\mathbf{y}) 
= 
- 
\int 
p_y(\mathbf{y}) 
\ln p_y(\mathbf{y}) 
d\mathbf{y}.
\end{equation}
%
By the transformation
%
\begin{equation}
p_y(\mathbf{y}) 
= 
p_x(\mathbf{A} \mathbf{x}) 
\abs{\det \mathbf{A} ^ {-1}},
\end{equation}
%
the right hand side can be written as
%
\begin{equation}
- 
\int 
p_x(\mathbf{A} \mathbf{x}) 
\ln p_x(\mathbf{A} \mathbf{x}) 
\abs{\det \mathbf{A}} 
d\mathbf{x} 
- 
\ln \abs{\det \mathbf{A} ^ {-1}} 
\int 
p_y(\mathbf{y}) 
d\mathbf{y}.
\end{equation}
%
By the transformation 
%
\begin{equation}
\mathbf{x}' 
= 
\mathbf{A} 
\mathbf{x},
\end{equation}
%
the first term can be written as
%
\begin{equation}
- 
\int 
p_x(\mathbf{x}') 
\ln p_x(\mathbf{x}') 
d\mathbf{x}' 
= 
{\rm H}(\mathbf{x}),
\end{equation}
%
and the second term can be written as
%
\begin{equation}
- 
\ln 
\abs{\det \mathbf{A} ^ {-1}} 
= \ln \abs{\det \mathbf{A}}.
\end{equation}
%
Therefore,
%
\begin{equation}
{\rm H}(\mathbf{y}) 
= 
{\rm H}(\mathbf{x}) 
+ 
\ln \abs{\det \mathbf{A}}.
\end{equation}
%


\subsection{}
\label{1.33}
Let $x$ and $y$ be two discrete variables.
By the definition,
%
\begin{equation}
{\rm H} (y | x) 
= 
- 
\sum_{i} 
\sum_{j} 
p(x_i, y_j) 
\ln p(y_j | x_i).
\end{equation}
%
If ${\rm H}(y | x)$ is zero, then
%
\begin{equation}
0 
= 
- 
\sum_{i} 
p(x_i)
\sum_{j} 
p(y_j | x_i) 
\ln p(y_j | x_i).
\end{equation}
%
Since
%
\begin{equation}
\begin{aligned}
p(x_i) 
&\geq 
0, 
\\
p(y_j | x_i) 
\ln p(y_j | x_i) 
&\leq 
0.
\end{aligned}
\end{equation}
%
for all $i$ and $j$, the equation reduces to
%
\begin{equation}
p(y_j | x_i) 
\ln p(y_j | x_i) 
= 
0.
\end{equation}
%
Therefore, $p(y_j | x_i)$ is zero or one.
Thus, since
%
\begin{equation}
\sum_j p(y_j | x_i) 
= 
1,
\end{equation}
%
$p(y_j | x_i)$ is one for a unique $x_i$ and zero for others.


\subsection{}
\label{1.34}
Let
%
\begin{equation}
\begin{aligned}
L(p(x)) 
=&
- 
\int_{- \infty}^{\infty} 
p(x) 
\ln p(x) 
dx 
+ 
\lambda_1 
\left( 
	\int_{- \infty}^{\infty} 
	p(x) 
	dx 
	- 
	1 
\right) \\
&+ 
\lambda_2 
\left( 
	\int_{- \infty}^{\infty} 
	x 
	p(x) 
	dx 
	- 
	\mu 
\right) 
+ 
\lambda_3 
\left( 
	\int_{-\infty}^{\infty} 
	(x - \mu) ^ 2 
	p(x) 
	dx 
	- 
	\sigma ^ 2 
\right).
\end{aligned}
\end{equation}
%
Setting the derivtive with respect to $p(x)$ to zero gives
%
\begin{equation}
0 
= 
- 
\ln p(x) 
- 
1 
+ 
\lambda_1 
+ 
\lambda_2 
x 
+ 
\lambda_3 
(x - \mu) ^ 2.
\end{equation}
%
Therefore,
%
\begin{equation}
p(x) 
= 
\exp 
\left( 
	- 1 
	+ 
	\lambda_1 
	+ 
	\lambda_2 
	x 
	+ 
	\lambda_3 
	(x - \mu) ^ 2 
\right).
\end{equation}
%
Therefore,
%
\begin{equation}
p(x) 
= 
\exp 
\left( 
	- 
	1 
	+ 
	\lambda_1 
	- 
	\frac{\lambda_2 ^ 2}{4 \lambda_3} 
	+ 
	\lambda_3 
	\left( 
		x 
		- 
		\left( 
			\mu 
			- 
			\frac{\lambda_2}{2 \lambda_3} 
		\right) 
	\right) ^ 2 
\right).
\end{equation}
%
Substituting it to 
%
\begin{equation}
\begin{aligned}
\int_{- \infty}^{\infty} 
p(x) 
dx 
&= 
1, \\
\int_{- \infty}^{\infty} 
x 
p(x) 
dx 
&= 
\mu, \\
\int_{- \infty}^{\infty} 
(x - \mu) ^ 2 
p(x) 
dx 
&= 
\sigma ^ 2,
\end{aligned}
\end{equation}
%
gives
%
\begin{equation}
\begin{aligned}
\exp 
\left( 
	- 
	1 
	+ 
	\lambda_1 
	- 
	\frac{\lambda_2 ^ 2}{4 \lambda_3} 
\right) 
\int_{- \infty}^{\infty} 
\exp 
\left( 
	\lambda_3 
	\left( 
		x 
		- 
		\left( 
			\mu 
			- 
			\frac{\lambda_2}{2 \lambda_3} 
		\right) 
	\right) ^ 2 
\right) dx 
&= 
1, \\
\exp 
\left( 
	- 
	1 
	+ 
	\lambda_1 
	- 
	\frac{\lambda_2 ^ 2}{4 \lambda_3} 
\right) 
\int_{- \infty}^{\infty} 
x 
\exp 
\left( 
	\lambda_3 
	\left( 
		x 
		- 
		\left( 
			\mu 
			- 
			\frac{\lambda_2}{2 \lambda_3} 
		\right) 
	\right) ^ 2 
\right) 
dx 
&= 
\mu, \\
\exp 
\left( 
	- 
	1 
	+ 
	\lambda_1 
	- 
	\frac{\lambda_2 ^ 2}{4 \lambda_3} 
\right) 
\int_{- \infty}^{\infty} 
(x - \mu) ^ 2 
\exp 
\left( 
	\lambda_3 
	\left( 
		x 
		- 
		\left( 
			\mu 
			-
			 \frac{\lambda_2}{2 \lambda_3} 
		\right) 
	\right) ^ 2 
\right) 
dx 
&= 
\sigma ^ 2.
\end{aligned}
\end{equation}
%
By the transformation
%
\begin{equation}
y 
= 
\sqrt{- \lambda_3} 
\left( 
	x 
	- 
	\left( 
		\mu 
		- 
		\frac{\lambda_2}{2 \lambda_3} 
	\right) 
\right),
\end{equation}
%
they can be written as
%
\begin{equation}
\begin{aligned}
\exp 
\left( 
	- 
	1 
	+ 
	\lambda_1 
	- 
	\frac{\lambda_2 ^ 2}{4 \lambda_3} 
\right) 
\int_{- \infty}^{\infty} 
\exp 
\left( 
	- 
	y ^ 2 
\right) 
(- \lambda_3) ^ {- \frac{1}{2}} 
dy 
&= 
1,\\
\exp 
\left( 
	- 
	1 
	+ 
	\lambda_1 
	- 
	\frac{\lambda_2 ^ 2}{4 \lambda_3} 
\right) 
\int_{- \infty}^{\infty} 
\left( 
	(- \lambda_3) ^ {- \frac{1}{2}} 
	y 
	+ 
	\mu 
	- 
	\frac{\lambda_2}{2 \lambda_3} 
\right) 
\exp 
\left( 
	- 
	y ^ 2 
\right) 
(- \lambda_3) ^ {- \frac{1}{2}} 
dy 
&= 
\mu,\\
\exp 
\left( 
	- 
	1 
	+ 
	\lambda_1 
	- 
	\frac{\lambda_2 ^ 2}{4 \lambda_3} 
\right) 
\int_{- \infty}^{\infty} 
\left( 
	(- \lambda_3) ^ {- \frac{1}{2}} 
	y 
	- 
	\frac{\lambda_2}{2 \lambda_3} 
\right) ^ 2 
\exp 
\left( 
	- 
	y ^ 2 
\right) 
(- \lambda_3) ^ {- \frac{1}{2}} 
dy 
&= 
\sigma ^ 2.
\end{aligned}
\end{equation}
%
Since
%
\begin{equation}
\begin{aligned}
\int_{- \infty}^{\infty} 
\exp 
\left( 
	- 
	y ^ 2 
\right) 
dy 
&= 
\Gamma 
\left( 
	\frac{1}{2} 
\right), \\
\int_{- \infty}^{\infty} 
y 
\exp 
\left( 
	- 
	y ^ 2 
\right) 
dy 
&= 
0, \\
\int_{- \infty}^{\infty} 
y ^ 2 
\exp 
\left( 
	- 
	y ^ 2 
\right) 
dy 
&= 
\Gamma 
\left( 
	\frac{3}{2} 
\right),
\end{aligned}
\end{equation}
%
they can be written as
%
\begin{equation}
\begin{aligned}
\exp 
\left( 
	- 
	1 
	+ 
	\lambda_1 
	- 
	\frac{\lambda_2 ^ 2}{4 \lambda_3} 
\right) 
(- \lambda_3) ^ {- \frac{1}{2}} 
\Gamma 
\left( 
	\frac{1}{2} 
\right) 
&= 
1, \\
\exp 
\left( 
	- 
	1 
	+ 
	\lambda_1 
	- 
	\frac{\lambda_2 ^ 2}{4 \lambda_3} 
\right) 
\left( 
	\mu 
	- 
	\frac{\lambda_2}{2 \lambda_3} 
\right) 
(- \lambda_3) ^ {- \frac{1}{2}} 
\Gamma 
\left( 
	\frac{1}{2} 
\right) &= 
\mu, \\
\exp 
\left( 
	- 
	1 
	+ 
	\lambda_1 
	- 
	\frac{\lambda_2 ^ 2}{4 \lambda_3} 
\right) 
\left( 
	(- \lambda_3) ^ {- \frac{3}{2}} 
	\Gamma 
	\left( 
		\frac{3}{2} 
	\right) 
	+ 
	(- \lambda_3) ^ {- \frac{1}{2}} 
	\frac{\lambda_2 ^ 2}{4 \lambda_3 ^ 2} 
	\Gamma 	
	\left( 
		\frac{1}{2} 
	\right) 
\right) 
&= 
\sigma ^ 2.
\end{aligned}
\end{equation}
%
Therefore,
%
\begin{equation}
\begin{aligned}
\lambda_1 
&=
1 
- 
\frac{1}{2} 
\ln 
\left( 
	2 \pi \sigma ^ 2 
\right), \\
\lambda_2 
&= 
0, \\
\lambda_3 
&= 
- \frac{1}{2 \sigma ^ 2}.
\end{aligned}
\end{equation}
%
Thus,
%
\begin{equation}
p(x) 
= 
\left( 
	2 \pi \sigma ^ 2 
\right) ^ {- \frac{1}{2}} 
\exp 
\left( 
	- 
	\frac{1}{2 \sigma ^ 2} 
	(x - \mu) ^ 2 
\right).
\end{equation}
%


\subsection{}
\label{1.35}
Let $x$ be a variable such that
%
\begin{equation}
p(x) 
= 
\mathcal{N} 
\left(
	x 
	| 
	\mu, 
	\sigma ^ 2
\right).
\end{equation}
%
Then, by the definition,
%
\begin{equation}
{\rm H}(x) 
= 
- 
\int_{- \infty}^{\infty} 
\mathcal{N} 
\left( 
	x 
	| 
	\mu, 
	\sigma ^ 2 
\right) 
\ln \mathcal{N} 
\left( 
	x 
	| 
	\mu, 
	\sigma ^ 2 
\right) 
dx.
\end{equation}
%
The right hand side can be written as
%
\begin{equation}
\begin{aligned}
&
- 
\int_{- \infty}^{\infty} 
\mathcal{N} 
\left( 
	x 
	| 
	\mu, 
	\sigma ^ 2 
\right) 
\left( 
	- 
	\frac{1}{2} 
	\ln 
	\left( 
		2 \pi \sigma ^ 2 
	\right) 
	- 
	\frac{1}{2 \sigma ^2} 
	(x - \mu) ^ 2 
\right) 
dx \\
&= 
\frac{1}{2} 
\ln 
\left( 
	2 \pi \sigma ^ 2 
\right) 
\int_{- \infty}^{\infty} 
\mathcal{N} 
\left( 
	x 
	| 
	\mu, 
	\sigma ^ 2 
\right) 
dx 
+ 
\frac{1}{2 \sigma ^2} 
\int_{- \infty}^{\infty} 
(x - \mu) ^ 2 
\mathcal{N} 
\left( 
	x 
	| 
	\mu, 
	\sigma ^ 2 
\right) 
dx.
\end{aligned}
\end{equation}
%
Therefore,
%
\begin{equation}
{\rm H}(x) 
= 
\frac{1}{2} 
\left( 
	1 
	+ 
	\ln 
	\left( 
		2 \pi \sigma ^ 2 
	\right) 
\right).
\end{equation}
%


\subsection{(Incomplete)}
\label{1.36}
Let $f$ be a strictly convex function.
Then, by the definition,
%
\begin{equation}
f 
\left( 
	\lambda 
	a 
	+ 
	(1 - \lambda) 
	b 
\right) 
\leq 
\lambda 
f(a)
+ 
(1 - \lambda) 
f(b),
\end{equation}
%
where $a \leq b$ and $0 \leq \lambda \leq 1$.
Let
%
\begin{equation}
x 
= 
\lambda 
a 
+ 
(1 - \lambda) 
b.
\end{equation}
%
Then, the inequality can be written as
%
\begin{equation}
f(x) 
\leq 
\frac{b - x}{b - a} 
f(a) 
+ 
\frac{x - a}{b - a} 
f(b).
\end{equation}
%
Let 
%
\begin{equation}
g(x) 
= 
\frac{b - x}{b - a} 
f(a) 
+ 
\frac{x - a}{b - a} 
f(b)
- 
f(x).
\end{equation}
%
Then,
%
\begin{equation}
g(x) 
\geq 
0.
\end{equation}
%
Additionally, for $x > a$,
%
\begin{equation}
g(x) 
= 
(x - a) 
\left( 
	\frac{f(b) - f(a)}{b - a} 
	- 
	\frac{f(x) - f(a)}{x - a} 
\right).
\end{equation}
%
By the mean value theorem, there exists $c$ and $y$ such that $a \leq c \leq b$, $a \leq y \leq x$ and
%
\begin{equation}
\begin{aligned}
f'(c) 
&= 
\frac{f(b) - f(a)}{b - a}, 
\\
f'(y) 
&= 
\frac{f(x) - f(a)}{x - a}.
\end{aligned}
\end{equation}
%
Then, for $x > a$, the inequality reduces to
%
\begin{equation}
f'(y) 
\leq 
f'(c). 
\end{equation}
%



\subsection{}
\label{1.37}
Let $\mathbf{x}$ and $\mathbf{y}$ be two variables.
Then, by the definition, 
%
\begin{equation}
{\rm H}(\mathbf{x}, \mathbf{y}) 
= 
- 
\int 
\int 
p(\mathbf{x}, \mathbf{y}) 
\ln p(\mathbf{x}, \mathbf{y}) 
d\mathbf{x} 
d\mathbf{y}.
\end{equation}
%
The right hand side can be written as
%
\begin{equation}
\begin{aligned}
&
- 
\int 
\int 
p(\mathbf{x}, \mathbf{y}) 
\left( 
	\ln p(\mathbf{y} | \mathbf{x}) 
	+ 
	\ln p(\mathbf{x}) 
\right) 
d\mathbf{x} 
d\mathbf{y} \\
&= 
- 
\int 
\int 
p(\mathbf{x}, \mathbf{y}) 
\ln p(\mathbf{y} | \mathbf{x})
d\mathbf{x} 
d\mathbf{y} 
- 
\int 
\left( 
	\int p(\mathbf{x}, \mathbf{y}) 
	d\mathbf{y} 
\right) 
\ln p(\mathbf{x}) 
d\mathbf{x}.
\end{aligned}
\end{equation}
%
By the definition, the first term of the right hand side can be written as ${\rm H}(\mathbf{y} | \mathbf{x})$ and the second term can be written as ${\rm H}(\mathbf{x})$.
Therefore,
%
\begin{equation}
{\rm H}(\mathbf{x}, \mathbf{y}) 
= 
{\rm H}(\mathbf{y} | \mathbf{x}) 
+ 
{\rm H}(\mathbf{x}).
\end{equation}
%


\subsection{}
\label{1.38}
Let $f$ be a strictly convex function.
Then, by the definition,
%
\begin{equation}
f 
\left( 
	\lambda 
	x_1 
	+ 
	(1 - \lambda) 
	x_2 
\right) 
\leq 
\lambda 
f(x_1) 
+ 
(1 - \lambda) 
f(x_2),
\end{equation}
%
where $0 \leq \lambda \leq 1$.
Let us assume that 
%
\begin{equation}
f 
\left( 
	\sum_{i = 1}^{M} 
	\lambda_i 
	x_i 
\right) 
\leq 
\sum_{i = 1}^{M} 
\lambda_i 
f(x_i),
\end{equation}
%
where $\lambda_i \geq 0$ and
%
\begin{equation}
\sum_{i = 1}^{M} 
\lambda_i 
= 
1.
\end{equation}
%
Here, let $\lambda_i \geq 0$ and
%
\begin{equation}
\sum_{i = 1}^{M + 1} 
\lambda_i 
= 
1.
\end{equation}
%
Then, by the definition,
%
\begin{equation}
f 
\left( 
	\sum_{i = 1}^{M + 1} 
	\lambda_i x_i 
\right) 
\leq 
\lambda_{M + 1} 
f ( x_{M + 1} ) 
+ 
(1 - \lambda_{M + 1}) 
f 
\left( 
	\sum_{i = 1}^{M} 
	\frac{\lambda_i}{1 - \lambda_{M + 1}} 
	x_i 
\right).
\end{equation}
%
By the assumption,
%
\begin{equation}
f 
\left( 
	\sum_{i = 1}^{M} 
	\frac{\lambda_i}{1 - \lambda_{M + 1}} 
	x_i 
\right) 
\leq 
\sum_{i = 1}^{M} 
\frac{\lambda_i}{1 - \lambda_{M + 1}} 
f(x_i).
\end{equation}
%
Therefore,
%
\begin{equation}
f 
\left( 
	\sum_{i = 1}^{M + 1} 
	\lambda_i 
	x_i 
\right) 
\leq 
\lambda_{M + 1} 
f ( x_{M + 1} ) 
+ 
(1 - \lambda_{M + 1}) 
\sum_{i = 1}^{M} 
\frac{\lambda_i}{1 - \lambda_{M + 1}} 
f(x_i).
\end{equation}
%
Thus,
%
\begin{equation}
f 
\left( 
	\sum_{i = 1}^{M + 1} 
	\lambda_i 
	x_i 
\right) 
\leq 
\sum_{i = 1}^{M + 1} 
\lambda_i 
f(x_i).
\end{equation}
%
Hence, the assumption is proved by induction on $M$.


\subsection{}
\label{1.39}
Let $x$ and $y$ be two binary variables where
%
\begin{equation}
\begin{aligned}
p(x = 0, y = 0) 
&= 
\frac{1}{3}, \\
p(x = 0, y = 1) 
&= 
\frac{1}{3}, \\
p(x = 1, y = 0) 
&= 
0, \\
p(x = 1, y = 1) 
&= 
\frac{1}{3}.
\end{aligned}
\end{equation}
%


\subsubsection{}
By the definition,
%
\begin{equation}
{\rm H}(x) 
= 
- 
\sum 
p(x)
\ln p(x).
\end{equation}
%
By the distribution,
%
\begin{equation}
\begin{aligned}
p(x = 0) 
&= 
\frac{2}{3}, \\
p(x = 1) 
&= 
\frac{1}{3}.
\end{aligned}
\end{equation}
%
Therefore,
%
\begin{equation}
{\rm H}(x) 
= 
\ln 3 
- 
\frac{2}{3} 
\ln 2.
\end{equation}
%


\subsubsection{}
By the definition,
%
\begin{equation}
{\rm H}(y) 
= 
- 
\sum 
p(y) 
\ln p(y).
\end{equation}
%
By the distribution,
%
\begin{equation}
\begin{aligned}
p(y = 0) 
&= 
\frac{1}{3}, 
\\
p(y = 1) 
&= 
\frac{2}{3}.
\end{aligned}
\end{equation}
%
Therefore,
%
\begin{equation}
{\rm H}(y) 
= 
\ln 3 
- 
\frac{2}{3} \ln 2.
\end{equation}
%


\subsubsection{}
By the definition,
%
\begin{equation}
{\rm H}(y | x) 
= 
- 
\sum 
p(x, y) 
\ln p(y | x).
\end{equation}
%
By the definition,
%
\begin{equation}
\begin{aligned}
p(y = 0 | x = 0) 
&= 
\frac{p(x = 0, y = 0)}{p(x = 0)}, 
\\
p(y = 0 | x = 1) 
&= 
\frac{p(x = 1, y = 0)}{p(x = 1)}, 
\\
p(y = 1 | x = 0) 
&= 
\frac{p(x = 0, y = 1)}{p(x = 0)}, 
\\
p(y = 1 | x = 1) 
&= 
\frac{p(x = 1, y = 1)}{p(x = 1)}.
\end{aligned}
\end{equation}
%
Then, by the distribution,
%
\begin{equation}
\begin{aligned}
p(y = 0 | x = 0) 
&=
\frac{1}{2}, 
\\
p(y = 0 | x = 1) 
&= 
0, 
\\
p(y = 1 | x = 0) 
&= 
\frac{1}{2}, 
\\
p(y = 1 | x = 1) 
&= 
1.
\end{aligned}
\end{equation}
%
Therefore,
%
\begin{equation}
{\rm H}(y | x) 
= 
\frac{2}{3} 
\ln 2.
\end{equation}
%


\subsubsection{}
By the definition,
%
\begin{equation}
{\rm H}(x | y) 
= 
- 
\sum 
p(x, y) 
\ln p(x | y).
\end{equation}
%
By the definition,
%
\begin{equation}
\begin{aligned}
p(x = 0 | y = 0) 
&= 
\frac{p(x = 0, y = 0)}{p(y = 0)}, 
\\
p(x = 0 | y = 1) 
&= 
\frac{p(x = 0, y = 1)}{p(y = 1)}, 
\\
p(x = 1 | y = 0) 
&= 
\frac{p(x = 1, y = 0)}{p(y = 0)}, 
\\
p(x = 1 | y = 1) 
&= 
\frac{p(x = 1, y = 1)}{p(y = 1)}.
\end{aligned}
\end{equation}
%
Then, by the distribution,
%
\begin{equation}
\begin{aligned}
p(x = 0 | y = 0) 
&= 
1, 
\\
p(x = 0 | y = 1) 
&= 
\frac{1}{2}, 
\\
p(x = 1 | y = 0) 
&= 
0, 
\\
p(x = 1 | y = 1) 
&= 
\frac{1}{2}.
\end{aligned}
\end{equation}
%
Therefore,
%
\begin{equation}
{\rm H}(x | y) 
= 
\frac{2}{3} 
\ln 2.
\end{equation}
%


\subsubsection{}
By the definition,
%
\begin{equation}
{\rm H}(x, y) 
= 
- 
\sum 
p(x, y) 
\ln p(x, y).
\end{equation}
%
Therefore,
%
\begin{equation}
{\rm H}(x, y) 
= 
\ln 3.
\end{equation}
%


\subsubsection{}
By the definition,
%
\begin{equation}
{\rm I}(x, y) 
= 
- 
\sum p(x, y) 
\ln \frac{p(x) p(y)}{p(x, y)}.
\end{equation}
%
By the distribution, the right hand side can be written as
%
\begin{equation}
{\rm H}(x) 
+ 
{\rm H}(y) 
- 
{\rm H}(x, y).
\end{equation}
%
Therefore,
%
\begin{equation}
{\rm I}(x, y) 
= 
\ln 3
- 
\frac{4}{3} 
\ln 2.
\end{equation}
%


\subsection{}
\label{1.40}
Let $\{ x_i \}$ be a set of points where $x_i > 0$, and let $\{ \lambda_i \}$ be a set of coefficients where $\lambda_i \geq 0$ and 
%
\begin{equation}
\sum_{i = 1}^{M} 
\lambda_i 
= 
1.
\end{equation}
%
By Jensen's inequality,
%
\begin{equation}
\sum_{i = 1}^{M} 
\lambda_i 
\ln x_i 
\leq 
\ln 
\left( 
	\sum_{i = 1}^{M} 
	\lambda_i 
	x_i 
\right).
\end{equation}
%
Therefore,
%
\begin{equation}
\prod_{i = 1}^{M} 
x_i ^ {\lambda_i} 
\leq 
\sum_{i = 1}^{M} 
\lambda_i 
x_i.
\end{equation}
%
Substituting
%
\begin{equation}
\lambda_i 
= 
\frac{1}{M}
\end{equation}
%
gives
%
\begin{equation}
\left( 
	\prod_{i = 1}^{M} 
	x_i 
\right) ^ \frac{1}{M} 
\leq \
frac{1}{M} 
\sum_{i = 1}^{M} 
x_i.
\end{equation}
%


\subsection{}
\label{1.41}
Let $\mathbf{x}$ and $\mathbf{y}$ be continuous variables.
Then, by the definitnion,
%
\begin{equation}
{\rm I}(\mathbf{x}, \mathbf{y}) = - \int \int p(\mathbf{x}, \mathbf{y}) \ln \frac{p(\mathbf{x}) p(\mathbf{y})}{p(\mathbf{x}, \mathbf{y})} d\mathbf{x} d\mathbf{y}.
\end{equation}
%
The right hand side can be written as
%
\begin{equation}
\begin{aligned}
&- \int \int p(\mathbf{x}, \mathbf{y}) \left( \ln p(\mathbf{x}) + \ln \frac{p(\mathbf{y})}{p(\mathbf{x}, \mathbf{y})} \right) d\mathbf{x} d\mathbf{y} \\
&= - \int \left( \int p(\mathbf{x}, \mathbf{y}) d\mathbf{y} \right) \ln p(\mathbf{x}) d\mathbf{x} + \int \int p(\mathbf{x}, \mathbf{y}) \ln p(\mathbf{x} | \mathbf{y}) d\mathbf{x} d\mathbf{y}.
\end{aligned}
\end{equation}
%
By the definition, the first term of the right hand side can be written as ${\rm H}(\mathbf{x})$ and the second term can be written as $- {\rm H(\mathbf{x} | \mathbf{y})}$.
Therefore,
%
\begin{equation}
{\rm I}(\mathbf{x}, \mathbf{y}) = {\rm H}(\mathbf{x}) - {\rm H(\mathbf{x} | \mathbf{y})}.
\end{equation}
%
By the definition,
%
\begin{equation}
{\rm I}(\mathbf{x}, \mathbf{y}) = {\rm I}(\mathbf{y}, \mathbf{x}).
\end{equation}
%
Thus,
%
\begin{equation}
{\rm I}(\mathbf{x}, \mathbf{y}) = {\rm H}(\mathbf{y}) - {\rm H(\mathbf{y} | \mathbf{x})}.
\end{equation}
%

